\documentclass[11pt]{article}
\usepackage[T1]{fontenc}
\usepackage[utf8]{inputenc}
\usepackage{lmodern,amsmath,amssymb,amsthm,mathrsfs,mathtools,booktabs,longtable,array}
\usepackage[margin=1in]{geometry}
\usepackage[colorlinks=true,urlcolor=blue,linkcolor=blue]{hyperref}
\newtheorem{theorem}{Theorem}[section]
\newtheorem{proposition}[theorem]{Proposition}
\newtheorem{lemma}[theorem]{Lemma}
\newtheorem{corollary}[theorem]{Corollary}
\theoremstyle{remark}\newtheorem{remark}[theorem]{Remark}
\DeclareMathOperator{\CH}{CH}
\DeclareMathOperator{\Alg}{Alg}
\DeclareMathOperator{\Pic}{Pic}
\DeclareMathOperator{\Alb}{Alb}
\DeclareMathOperator{\NS}{NS}
\DeclareMathOperator{\End}{End}
\DeclareMathOperator{\Hom}{Hom}
\DeclareMathOperator{\Nm}{Nm}
\newcommand{\Q}{\mathbf Q}\newcommand{\Z}{\mathbf Z}\newcommand{\F}{\mathbf F}
\newcommand{\tun}[1]{#1}
\setlength{\emergencystretch}{2em}
\title{C44\_POS: primitive signs, special axes, and the 1964 programme}
\author{Research continuation of JAC1964 and SIGN\_NEXT}
\date{22 November 2026}
\begin{document}
\maketitle
\begin{abstract}
Working research manuscript. The general fixed-axis sign and the Fermat
Abel--Jacobi injection remain unresolved. New checked results include
an exact axis primitive decomposition, lifts of every embedded smooth
axis (without a cycle-lifting assertion), explicit positive axis
subspaces, special-cover generic-stratum calculations, exact integral
arithmetic obstructions, a characteristic-$p$ detector/division theorem,
and geometric extensions of the primitive incidence polarization.
The proposed supersingular Fermat test case already satisfies the
required Hodge standard conjecture; its corrected restriction and
weighted Gysin conditions are verified explicitly.
Every use of 032 or Milne's 2026 report remains conditional on those
unrefereed preprints. The separate Milne 2022 transfer assertion is
exposed wherever used. Whole-proof reading events and finite
computations are distinguished in \texttt{CHECK\_EVENTS.tsv}.
\end{abstract}
\tableofcontents
\part{Checked findings and explicitly conditional deductions}
\section{Normalization, scope, and inherited results}
The starting isometry is JAC1964 H9-T7, equation (H9.18):
\[
 A_X\oplus A_B\simeq\CH^m(Y)_{\mathrm{hom},\Q}/\mathcal H_0,
 \qquad Q((z,w),(z',w'))=(-1)^m\int_Xzz'+(-1)^{m-1}\int_Bww'.
\]
Here $A_B=\ker\kappa_*$ on the algebraic middle classes of the axis,
and $h_BA_B=0$. This is an inherited statement, not a new result of
this run. The ambient sign is conditional on 032 Corollary 8.4 and
Milne's report. Weak semidefiniteness and strict positive definiteness
are distinguished throughout; no numerical/homological equivalence
on an arbitrary axis is silently assumed.

Geometric generic parameter points are taken over algebraic closures
of function fields. They are not closed points over
$\overline{\mathbf F}_p$. A monodromy vanishing assertion at the
former does not settle a sign at every latter point.

Every finding retains its stable C44 identifier. A displayed
proposition that assumes an additional unproved input is a checked
conditional implication, not a proof of that input. Prior statements
retain their JAC1964 or SIGN\_NEXT identifiers and are not claimed
as findings of this continuation. The governing source is the attached
\texttt{PROMPT(20261121-231316).md}.

% BEGIN work/axis/axis_contribution.tex
% C44_POS /root/axis_independent; GPT-6 Astra Pro; effort setting not exposed.
% Mathematical deductions for root review; no claim of formal verification.
\section{The exact axis primitive space and the limits of the proposed transfers}

\paragraph{Notation and carried input.}
Put $r=m-1$, so that $\dim X=2r+2$ and $\dim B=2r$.
Write $L_X$ and $L_B$ for multiplication by $h$ and $h_B$.
All cycle groups have rational coefficients, and $\operatorname{Alg}^j$
means the image of the rational cycle-class map for one fixed
$\ell\ne p$. No equality between numerical and homological equivalence
on $B$ is assumed. The carried statement JAC1964 H9-T7, equation
(H9.18), is the isometry
\[
 A_X\oplus A_B\ \simeq\ \mathcal H/\mathcal H_0,
 \qquad
 Q((z,w),(z',w'))=(-1)^m\int_Xzz'+(-1)^r\int_Bww'.
\]
Its proof establishes $L_Bw=0$ for $w\in A_B$, not merely
$L_B^2w=0$. The factorization $B\hookrightarrow Y\hookrightarrow X$
through a smooth pencil member and the injectivity of the latter
Gysin map are essential to that stronger conclusion. We use this
checked input without reproving its height isometry.

\subsection*{C44-A01. The axis gate is the entire remaining primitive sign on the axis}

\textbf{Proposition.}
There is an orthogonal direct sum of rational algebraic classes
\[
 \operatorname{Alg}^{r}(B)\cap P^{2r}(B)
 =\kappa^*\bigl(\operatorname{Alg}^{r}(X)\cap P^{2r}(X)\bigr)
       \mathbin{\perp} A_B,                                  \tag{C44.A01.1}
\]
where $P^{2r}(B)=\ker L_B$ and $P^{2r}(X)=\ker L_X^3$.
On the first summand the signed middle form is
\[
   (-1)^r\int_B\kappa^*a\,\kappa^*a'
       =(-1)^r\int_Xh^2aa'.                                 \tag{C44.A01.2}
\]
For every $0\le j<r$, restriction identifies the primitive algebraic
classes of codimension $j$ on $X$ and $B$, including their signed
Lefschetz pairings.
Consequently, \emph{conditional on the unrefereed 032/Milne input
asserting the Hodge standard conjecture for $X$}, strict positivity
on $A_B$ is equivalent to the full Hodge standard conjecture $I(B)$
for $h_B$. Weak nonnegativity on $A_B$ is equivalent to the version
of $I(B)$ with every positive-definite assertion replaced by
positive-semidefinite. These last two statements must not be
identified unless nondegeneracy on algebraic classes is additionally
known.

\textbf{Proof.}
Hard Lefschetz makes $L_X^2:H^{2r}(X)\longrightarrow H^{2r+4}(X)$
an isomorphism. Its inverse $\beta$ is induced by an algebraic
correspondence because the Lefschetz standard conjecture is known
for abelian varieties, without 032. The self-intersection and
projection formulas give
$\kappa_*\kappa^*=L_X^2$. Hence
$e=\kappa^*\beta\kappa_*$ is a projector and every algebraic middle
class $x$ of $B$ has a unique decomposition
\[
 x=\kappa^*a+w,\qquad a=\beta\kappa_*x,\qquad w\in A_B.       \tag{C44.A01.3}
\]
Both components are rational algebraic classes. The summands are
orthogonal because
$\int_B\kappa^*a\,w=\int_Xa\,\kappa_*w=0$.
H9-T7 gives $L_Bw=0$. The Gysin map
\[
 \kappa_*:H^{2r+2}(B)\longrightarrow H^{2r+6}(X)
\]
is an isomorphism, being dual to weak Lefschetz in degree $2r-2$.
Therefore $L_B\kappa^*a=0$ if and only if
$\kappa_*L_B\kappa^*a=L_X^3a=0$. This proves (C44.A01.1).
The projection formula proves (C44.A01.2).

For $j<r$, weak Lefschetz gives an isomorphism
$\kappa^*:H^{2j}(X)\longrightarrow H^{2j}(B)$.
Its inverse on this degree is the algebraic operator
\[
 (L_X^{2r+2-2j})^{-1}\kappa_*L_B^{2r-2j};                   \tag{C44.A01.4}
\]
indeed its composition with $\kappa^*$ is the identity.
Thus this is an isomorphism on rational algebraic classes too.
For $1\le j<r$, the map $\kappa_*$ in the degree of
$L_B^{2r-2j+1}\kappa^*a$ is an isomorphism by dual weak
Lefschetz in degree $2j-2$. Thus the primitive condition on $B$
is equivalent to
\[
 L_X^{2r-2j+3}a=0,
\]
which is precisely the primitive condition on $X$. For $j=0$
both primitive spaces are $H^0$, so the same assertion holds.
Finally
\[
 (-1)^j\int_B h_B^{2r-2j}\kappa^*a\,\kappa^*a'
 =(-1)^j\int_Xh^{2r+2-2j}aa'.
\]
The conditional positivity assertions now follow from the
orthogonal direct sum and the signed Hodge standard conjecture on
$X$. No positivity assertion on the new summand has been proved.
\hfill$\square$

\paragraph{Programme location and scope.}
This is a sharpened identification of JAC1964 QH9-axis, and concerns
cote 44, Th.~5, scan 004, with the primitive margin on scan 008.
In cote 16 terminology it isolates $I(B)$; the algebraic operator
$\beta$ uses $C_G(X)=B(X)$, not $I(B)$. It does not close
QH9-axis and does not upgrade its requested semidefinite sign to
strict positivity.

\subsection*{C44-A02. Every graph transfer to an abelian target annihilates the axis classes}

\textbf{Proposition.}
Let $D$ be an abelian variety and $f:B\longrightarrow D$ any
morphism. Then $f_*w=0$ in cohomology for every $w\in A_B$.
The same holds for a rational map, interpreted through its unique
extension. Consequently $A_B$ is orthogonal to every cohomology
class pulled back by such a morphism in the complementary degree.
This is unconditional. It does not rule out correspondences on
$B\times D$ which are not assembled from these graph maps.

\textbf{Proof.}
Choose a geometric base point $b$ on $B$. Let
$a_B:B\longrightarrow\operatorname{Alb}(B)$ be the pointed Albanese
map and let $u:\operatorname{Alb}(B)\longrightarrow X$ be induced by
$x\mapsto\kappa(x)-\kappa(b)$.
Weak Lefschetz in degree one identifies $H^1(X)$ with $H^1(B)$.
The identification of degree-one cohomology with the dual of the
Albanese Tate module therefore makes $u$ an isogeny.
After subtracting the constant $f(b)$, the morphism $f$ factors as
$v a_B$. Since $u$ is invertible in the isogeny category, there are
an integer $N>0$ and a homomorphism $t:X\longrightarrow D$ such that
$t u=[N]_D v$. Thus
\[
 [N]_D\circ f=t\circ\kappa+\text{a constant}.
\]
Translations on an abelian variety act trivially on cohomology.
Functoriality of proper pushforward gives
\[
 [N]_{D,*}f_*w=t_*\kappa_*w=0.
\]
Multiplication by $N$ induces an invertible map on rational
$\ell$-adic cohomology, so $f_*w=0$. The orthogonality assertion is
projection formula. A rational map from a smooth variety to an
abelian variety extends by Weil's extension theorem; one exact
modern locator is Milne, \emph{Abelian Varieties}, course notes
version 2.00, I, Theorem 3.2, printed p.~17 (PDF p.~23).\hfill$\square$

\paragraph{Programme location and scope.}
This specifies a limit of the proposed curve-test attack on
QH9-axis. It concerns the incidence higher Picard framework of
cote 44, Th.~1--4, scan 002, and the incidence divisor on scans
008--009. It complements JAC-P19 (no abelian rational domination)
and the already proved annihilation calculation in H9-T13, rather
than replacing either. An algebraic correspondence shifts
cohomological degree by an even number, so $H^1(C)$ by itself
cannot map into $H^{2r}(B)$ through an algebraic correspondence.
The elliptic tensor construction $w\boxtimes P_E$ resolves that
parity issue but its ambient pushforward is zero, exactly as
recorded in H9-T13. The missing input is a genuinely new
pairing-compatible correspondence, not a graph of another
Albanese map.

\subsection*{C44-A03. The finite-field test is absence of potential Tate roots, not one exceptional slope}

\textbf{Proposition.}
Suppose the axis and the embedding are defined over $\mathbf F_q$.
Let
\[
 V_B=\ker\bigl(\kappa_*:H^{2r}(B_{\overline{\mathbf F}_q},
 \mathbf Q_\ell)\longrightarrow H^{2r+4}(X_{\overline{\mathbf F}_q},
 \mathbf Q_\ell)(2)\bigr)
\]
be the untwisted primitive variable cohomology, with the twists in
the target understood by Gysin compatibility. Write $F$ for
geometric Frobenius and $P_B(t)=\det(t-F\mid V_B)$.
Every algebraic axis class is contained in
\[
 T_B=\bigcup_{a\ge1}\ker(F^a-q^{ra}\mid V_B).                \tag{C44.A03.1}
\]
In particular all its normalized Newton slopes are $r$.
If $P_B$ has no root $q^r\zeta$ for a root of unity $\zeta$,
then $A_B=0$.
For $D=\deg P_B>0$, a finite exact sufficient test is
\[
 \gcd\bigl(P_B(t),\ t^N-q^{rN}\bigr)=1,
 \qquad N=\operatorname{lcm}(1,2,\ldots,2D^2).               \tag{C44.A03.2}
\]
Here $P_B\in\mathbf Q[t]$ is obtained by dividing the middle
Frobenius polynomial of $B$ by that of the restricted ambient
summand. No Tate conjecture and no Frobenius semisimplicity
conjecture is needed for this vanishing criterion.

\textbf{Proof.}
A rational algebraic cycle on $B_{\overline{\mathbf F}_q}$ is
defined over a finite extension $\mathbf F_{q^a}$. Its codimension
$r$ cycle class in untwisted cohomology satisfies $F^a w=q^{ra}w$.
This proves (C44.A03.1). In crystalline realization, the same
cycle is a Frobenius eigenclass of slope $r$; equivalently every
root $q^r\zeta$ has normalized valuation $r$ at every prime over
$p$.
If the polynomial has no such root, every kernel in (C44.A03.1)
is zero, which proves the asserted vanishing.

The splitting (C44.A01.3) is Frobenius-compatible after the
polarization is defined over $\mathbf F_q$; its image is isomorphic
to $H^{2r}(X)$. Thus the usual rational Frobenius polynomials of
$X$ and $B$ give $P_B\in\mathbf Q[t]$ without asserting
rationality of a chosen algebraic basis of $A_B$.
If $q^r\zeta$ is a root and $n$ is the order of $\zeta$, the
cyclotomic polynomial of $\zeta$ divides
$q^{-rD}P_B(q^rt)$, so $\varphi(n)\le D$.
For $n=\prod p^a$,
\[
 \frac{\varphi(n)^2}{n}=\prod_{p^a\parallel n}p^{a-2}(p-1)^2
 \ge\frac12:
\]
each factor for $p\ge3$ is at least one, and the sole possible
factor less than one is $1/2$ for $2^1$. Consequently $n\le2D^2$,
so $n\mid N$. This proves (C44.A03.2), including its equivalence
to absence of roots of this form. For $D=0$ the conclusion is
immediate.\hfill$\square$

\paragraph{Why H9-T5 does not supply this test.}
JAC1964 H9-T5 is a theorem about the \emph{odd} variable local
system $R^{2m-1}$ of the pencil and its variable incidence abelian
factor. Its irreducibility hypothesis turns a nonzero injection
$H^1(D)(-(m-1))$ into an isomorphism; one slope outside
$[m-1,m]$ then contradicts that isomorphism. The even-degree
fixed-axis space in (C44.A03.1) is a different representation.
At a closed point the arithmetic Galois group is procyclic, and
an isolated algebraic class need not extend to a flat algebraic
class on the generic axis family. No family-monodromy
irreducibility statement can therefore be substituted into
(C44.A03.1) without a separate spreading theorem.
An out-of-range slope in some other direct summand is compatible
with a slope-$r$ Tate summand. The correct stronger sufficient
Newton condition here is absence of \emph{every} slope-$r$
subspace, while the cyclotomic criterion above is sharper.
This repairs the proposed use of H9-T5 without producing a
witness for every special axis.

\paragraph{Supersingular ambient varieties do not resolve this distinction.}
If $X$ is supersingular, $H^{2r}(X)$ has only slope $r$ because
$H^1(X)$ has slope $1/2$ and the cohomology algebra is exterior.
That identifies the slopes of the ambient summand in
(C44.A01.3), not the slopes or cycle classes of $V_B$.
A nonliftable algebraic class pulled back from $X$ lies in the
ambient summand and is removed by the projector $1-e$. Thus
ambient nonliftability alone gives no nonliftable element of
$A_B$.

\subsection*{C44-A05. The embedded axis itself always has a characteristic-zero lift}

\textbf{Proposition.}
Let $k=\overline{\mathbf F}_p$, let $X/k$ be an abelian variety,
let $L$ be ample, and let $s_0,s_1\in H^0(X,L)$ have a smooth
codimension-two common zero scheme $B$. There is a mixed-characteristic
discrete valuation ring $R$, with finite residue field contained in $k$,
an abelian scheme $\mathcal X/R$, a relatively ample line bundle
$\mathcal L$, and two sections whose common zero scheme
$\mathcal B/R$ is smooth and projective, such that the geometric
special fibre is the given $(X,L,s_0,s_1,B)$. In particular this
applies to every axis in QH9-axis, including supersingular ambient
abelian varieties. The proposition asserts no lifting of a specified
middle-dimensional algebraic class on $B$.

\textbf{Proof.}
Descend the given data to a finite subfield of $k$. The polarized
lifting theorem of Mumford--Norman--Oort gives an abelian scheme
with polarization lifting $(X,\lambda_L)$ over an integral
mixed-characteristic base. An exact statement is Oort,
\emph{Did Earlier Thoughts Inspire Grothendieck?}, Theorem (8.3).1,
p.~29. By finite-presentation descent we may use an integral base
$S$ of finite type over $\mathbf Z$, with a closed point of finite
residue field giving the prescribed geometric special fibre. The
scheme, group laws, and polarization all descend; passing to a
sufficiently large finite-type model preserves smoothness and
properness. This is the usual limit construction, as in Stacks
Project, Lemma 32.10.1, tag 01ZM, Lemma 32.8.9, tag 0C0C,
and Lemma 32.13.1, tag 081F.

The chosen line bundle also lifts after an \'etale neighborhood of
the special point: Olsson, \emph{Compactifications of moduli of
abelian varieties: an introduction}, Lemma 2.1.8 and its proof,
pp.~299--300. Its proof preserves the chosen special-fibre line
bundle. Alternatively, after choosing a representative of the
lifted polarization, the difference from $L$ lies in the dual
abelian variety and lifts after a smooth base change. Take an
integral component through the chosen point and still dominating
$S$; this is possible by flatness of the \'etale neighborhood.
Write $\pi:\mathcal X\to S$ for the resulting family and
$\mathcal L$ for its ample line bundle. For an ample line bundle
on an abelian scheme,
\[
 R^i\pi_*\mathcal L=0\quad(i>0),\qquad
 E=\pi_*\mathcal L\ \text{is locally free and commutes with base change}.
\]
The exact family statement is Olsson, \S2.1.4, p.~299.

The total space $W$ of the vector bundle $E\oplus E$ carries a
universal pair of sections on $\mathcal X\times_S W$. It is
integral of characteristic zero and of finite type over
$\mathbf Z$. The given pair determines a closed point $b\in W$
after a finite residue-field extension. Let $\mathcal Z$ denote
the universal common zero scheme. On $\mathcal Z$, failure of the
two relative differentials to have rank two is closed. Its image
in $W$ is closed because $\mathcal Z\to W$ is proper. The
complement $U$ contains $b$, since $B$ is smooth of codimension
two. The relative Jacobian criterion makes $\mathcal Z_U\to U$
smooth of relative dimension $\dim X-2$, and it is projective.
Thus $U$ is an integral finite-type mixed-characteristic
parameter scheme carrying a smooth projective lift of the
specified geometric fibre.

For completeness one can pass to an arithmetic trait without
changing that geometric fibre. Let $A=\mathcal O_{U,b}$, with
maximal ideal $\mathfrak m$, and put $d=\dim A$. This is an
excellent local domain, $0\ne p\in\mathfrak m$, and its residue
field is finite. Choose $x_1,\ldots,x_{d-1}$ so that
$(p,x_1,\ldots,x_{d-1})$ is $\mathfrak m$-primary. Set
$I=(x_1,\ldots,x_{d-1})$. Some minimal prime $P$ of $I$ does
not contain $p$: otherwise $p\in\sqrt I$ would give
$\sqrt I=\mathfrak m$, contrary to Krull's height theorem.
The local domain $A/P$ has dimension one, since $p$ is nonzero
and generates an ideal with maximal radical in that quotient.
Its normalization is finite by excellence; localization at a
point above its closed point is a discrete valuation ring $R$
with characteristic-zero fraction field and finite residue field.
Pulling back the family along
$\operatorname{Spec}R\to\operatorname{Spec}A\to U$ proves the
claim. The residue extension embeds in $k$, so the geometric
special fibre is the prescribed one.\hfill$\square$

\paragraph{Programme location and remaining obstruction.}
This completes the geometric lifting substep of QH9-axis(a),
in the schemes-over-$\mathbf Z$ framework. The unresolved issue
is surjectivity, for some such lift, of specialization on the
specified rational algebraic middle classes. Smooth proper base
change is an isomorphism on cohomology; it is not a lifting
theorem for cycles. The sign transfer to classes which do lift
is exactly the already settled H9-T8a specialization argument.
Cote 44, Th.~5, scan 004, and the primitive margin on scan 008
therefore retain the same cycle-level gate.

\section{Unresolved axis steps after the ordered attacks}

\paragraph{Lifting.}
For any specified smooth projective mixed-characteristic lift
of $(B,h_B)$ and any rational algebraic class on the geometric
generic fibre specializing to $w\in A_B$, compatibility of
cycle classes with smooth proper specialization carries
$L_Bw=0$ back to the same primitive relation on the generic
fibre. The specialization is an isomorphism on $\ell$-adic
cohomology, so a nonzero $w$ has nonzero generic class, and the
intersection number is unchanged. H9-T8a then gives
$(-1)^r\int_Bw^2>0$. This is the carried lifting argument,
not a surjectivity theorem for Chow specialization.

No proof that every $w\in A_B$ lifts with $B$ is obtained.
Nor is a specific cohomological class in $A_B$ proved
nonliftable in this contribution. Attempts using nonliftable
ambient Tate classes fail because their projection to $A_B$
is zero. Agugliaro, \emph{Examples for the standard conjecture
of Hodge type}, Doc. Math. 31 (2026), Proposition 5.12,
pp.~876--877, supplies a rank-two nonliftable Tate block for
CM lifts of certain abelian varieties; Remark 3.5, p.~870,
explains the passage to arbitrary abelian lifts. Even after
032 supplies algebraicity, these remain ambient examples.
They are not examples for the axis kernel.

\paragraph{Correspondences.}
C44-A02 and JAC1964 H9-T13 eliminate graph/ambient transfers.
A general correspondence from $B$ to an abelian variety could
still act nontrivially, but a splitting preserving the signed
geometric pairing has not been constructed. Merely splitting
an abstract Tate class from $h(B)$ would not identify that
pairing with a positive form.

\paragraph{Newton slopes.}
C44-A03 gives the correct finite-field vanishing test.
No Frobenius polynomial satisfying that test has been exhibited
for every special axis, and axes with $A_B\ne0$ necessarily
fail it. A slope witness in the odd pencil variation does not
settle the even fixed-axis sign.

\paragraph{Current gate.}
The unresolved assertion remains weak nonnegativity of
$(-1)^r\int_Bw^2$ for every algebraic $w\in\ker\kappa_*$.
Conditional on 032/Milne, C44-A01 identifies it with the weak
Hodge-standard assertion for $B$. No wrong-sign cycle is
constructed. No main question is marked closed by these
route boundaries.

% END work/axis/axis_contribution.tex

% BEGIN work/axis/root_explicit_axes.tex
\section{An explicit nonzero axis subspace in every characteristic}

\subsection*{C44-A04. Axes containing fibres of an abelian product}
\begin{theorem}[C44-A04]\label{thm:product-axis}
Let $r=m-1\ge2$, let $S$ and $T$ be abelian varieties of dimensions
$r$ and $r+2$, and put $X=S\times T$. Choose very ample classes
$s$ and $t$ on $S$ and $T$ and let $h=s+t$ on $X$. Let
$B\subset X$ be a smooth complete intersection of two divisors of
class $h$. Suppose that it contains the distinct fibres
$S_i=S\times\{t_i\}$ for $1\le i\le N$. Put
\[
 d_S=\int_Ss^r,\qquad d_T=\int_Tt^{r+2},\qquad
 a=\frac1{d_T}\sum_{j=0}^r(-1)^j(j+1)s^jt^{r-j},
 \qquad w_i=[S_i]_B-\kappa^*a.
 \tag{C44.A04.1}
\]
Then $w_i\in A_B$ and the exact signed Gram matrix is
\[
 \left((-1)^r\int_B w_iw_j\right)_{i,j}
       =(r+1)d_S\left(I_N-\frac1{d_T}{\bf1}{\bf1}^{t}\right).
 \tag{C44.A04.2}
\]
Smoothness of $B$ forces $N\le d_T$. Hence this matrix is positive
semidefinite in every characteristic, without 032. If $N<d_T$ it is
positive definite, and the $w_i$ are linearly independent rational
algebraic classes. In particular a single contained fibre supplies a
nonzero axis class of the required sign.

Smooth such axes containing a prescribed fibre exist after replacing
$s,t$ by sufficiently divisible positive multiples. The two defining
sections may be chosen to give a genuine Lefschetz pencil. Thus the
construction includes products of supersingular elliptic curves and
gives special axes with $A_B\ne0$ over $\overline{\F}_p$ for every $p$.
The assertion concerns the displayed subspace, not all of $A_B$.
\end{theorem}
\begin{proof}
All computations are in rational Weil cohomology, with their usual
cycle representatives. Since $s^{r+1}=0$ and $t^{r+3}=0$, the finite
identity $(1+z)^2\sum_{j=0}^r(-1)^j(j+1)z^j\equiv1\pmod{z^{r+1}}$
gives
\[
                      h^2a=t^{r+2}/d_T=[S_i]_X.
\]
Thus $\kappa_*w_i=0$ by self-intersection, and $w_i\in A_B$.
The normal bundle of $S_i$ in $X$ is trivial of rank $r+2$.
Because $B$ is smooth along $S_i$, its normal sequence is exact:
\[
 0\longrightarrow N_{S_i/B}\longrightarrow\mathcal O_S^{r+2}
   \longrightarrow L_S\oplus L_S\longrightarrow0.
\]
Here $c_1(L_S)=s$. Consequently
$c(N_{S_i/B})=(1+s)^{-2}$ and
\[
                 \int_B[S_i]^2=(-1)^r(r+1)d_S.
\]
Distinct fibres are disjoint, so their intersection is zero.
Restriction of $a$ to any $S_i$ has degree
$(-1)^r(r+1)d_S/d_T$; the same is true of
$\int_B(\kappa^*a)^2=\int_Xa h^2a$.
Expanding each $w_iw_j$ therefore gives (C44.A04.2).

We prove the bound on the number of contained fibres without using a
positivity conjecture. Write the two defining sections in a basis of
$H^0(S,L_S)$ as
$F_\nu=\sum_\alpha\sigma_\alpha\tau_{\nu\alpha}$,
where $\tau_{\nu\alpha}\in H^0(T,L_T)$ and $\nu=1,2$.
Let $Z\subset T$ be the common zero scheme of all coefficients
$\tau_{\nu\alpha}$. Its geometric points are exactly the fibres
contained in $B$. Suppose $T_zZ$ has positive dimension at a point
$z\in Z$. Differentiating the coefficients shows that the normal
surjection along $S\times\{z\}$ factors through
\[
       \mathcal O_S^{r+1}\longrightarrow L_S\oplus L_S
\]
(or through a trivial bundle of still smaller rank). Its kernel
would have rank at most $r-1$ but total Chern class $(1+s)^{-2}$.
Its $r$th Chern class would therefore be both zero, by rank, and
$(-1)^r(r+1)s^r$, whose degree is nonzero. This is impossible.
Thus $Z$ has zero tangent space at every geometric point; it is a
finite reduced scheme. Choose $r+2$ general linear combinations of
the coefficient sections. Since their common base scheme is finite,
successive dimension cuts yield a zero-dimensional complete
intersection containing $Z$. Its degree is $d_T$, so
$N\le\#Z\le d_T$.
The eigenvalues of the matrix in parentheses in (C44.A04.2) are
$1$ with multiplicity $N-1$ and $1-N/d_T$ with multiplicity one.
This proves all the asserted signs and, for $N<d_T$, independence.

For existence take the prescribed smooth centre $S_0=S\times\{0\}$.
For sufficiently positive $L_S,L_T$, the linear system
$H^0(X,I_{S_0}\otimes(L_S\boxtimes L_T))$ separates second jets
at each point off $S_0$, separates the first jets at each pair of
distinct points off $S_0$, and maps onto the conormal first jets
along $S_0$. These are uniform assertions for all such points:
apply relative Serre vanishing to the universal length-bounded
jet subschemes and the fixed coherent ideal $I_{S_0}$ (or use
high powers of product sections). A common sufficiently high
power satisfies all three conditions.
The latter assertion also follows directly from product sections:
sections of $L_S$ generate its fibres and sections of $I_0L_T$
generate the cotangent space of $T$ at $0$ after taking a high power.
For two general sections, the locus on $S_0$ where their normal
derivatives have rank at most one has expected codimension $r+1$;
the evaluation to the two normal covectors is surjective at each
point, so the incidence dimension estimate, with $\dim S_0=r$,
makes this locus empty. Away from $S_0$, jet separation and
Bertini give a smooth codimension-two intersection. Here are the
dimension counts for the pencil condition. Write $n=\dim X=2r+2$.
At a specified point off $S_0$, singularity imposes $n+1$
independent first-jet equations on a section. Degeneracy of its
quadratic part imposes an additional proper equation, by the
second-jet surjectivity. Letting the point vary gives a bad locus
of codimension at least two in the projective linear system.
Singularities at two specified distinct points impose $2(n+1)$
independent equations; their parameter space has dimension $2n$,
again giving codimension at least two. For singularity along the
imposed centre, the normal covector must vanish; this imposes
$r+2$ conditions at a specified point of $S_0$, and $\dim S_0=r$.
Its incidence image has codimension at least two as well. A general
projective line avoids these bad loci. Smoothness of its axis
also prevents the line from being tangent to the discriminant at
a single ordinary quadratic point: such tangency says that the
other pencil section vanishes at that point, putting a singular
member's zero differential into the axis's defining pair.
Thus the line gives a Lefschetz pencil, with distinct singular
values and no singularity on $S_0$. Since $n$ is even, the
nondegenerate quadratic forms are a nonempty open also in
characteristic two. These are the ordinary Lefschetz local models.
All conditions used are nonempty Zariski-open conditions over the
algebraically closed field. Over $\overline{\F}_p$ the chosen
sections are defined over a finite extension. This proves existence.
\end{proof}

\paragraph{C44-A04a: a concrete sixfold normalization.}
Take $X=E^6$, $S=E^2$, $T=E^4$ and
$h=d\sum_{i=1}^{6}[0]_i$, with $d$ sufficiently large for the
construction. For an axis containing $E^2\times\{0\}$,
\[
 d_S=2d^2,\qquad d_T=24d^4,\qquad
 \int_Bw^2=6d^2-\frac1{4d^2}>0.
\]
The integral cycle $\widetilde w=d_T[S]_B-
\kappa^*(t^2-2st+3s^2)$ has
\[
        \int_B\widetilde w^2=144d^6(24d^4-1)>0.
\]
By H9-T7 the corresponding height classes have these exact heights
in its normalization. One may choose $E$ supersingular in every
characteristic. These are positive special-axis classes, not a
classification of nonliftable classes and not a wrong-sign example.

\paragraph{Archive and source location.}
The height interpretation is JAC1964 H9-T7, (H9.18), for cote 44,
Th.~5, scan 004. The primitive qualification is the margin on scan
008. The normal-bundle and self-intersection calculations use Fulton,
\emph{Intersection Theory}, Section 6.3, self-intersection formula,
and the codimension-two blowup normalization already checked in
H9-T7. The existence argument uses the ordinary jet construction
underlying SGA~7~II, XVII, 1.1--1.2 (pp.~213--214), 2.5
(p.~217), 3.2.1--3.2.4 (p.~219), 3.3 (p.~221), 3.7
(p.~234), and 4.2--4.3 (pp.~236--239); its imposed-centre
incidence argument is provided here, not attributed as a separate
statement to SGA. In cote 16 language the calculation verifies $I$
on a specified algebraic primitive subspace. It does not establish
$I(B)$ on its complement or close QH9-axis.

% END work/axis/root_explicit_axes.tex

% BEGIN work/covers/contribution.tex
\section{Special cover strata and the old axis summand}

\subsection{C44-C01: the old variable classes do not disappear}
\begin{proposition}[C44-C01]\label{prop:old-axis-cover}
Use JAC1964 H9-T18. Write $A$ for the abelian fourfold, $B$ for its
surface axis, $X_0=\operatorname{Bl}_B A$, $M_0=A\times\mathbf P^1$,
and $k_0:X_0\hookrightarrow M_0$ for the incidence embedding.
For a finite separable cover $\pi:C\to\mathbf P^1$ of degree $b$,
unramified over the critical values, put
$X_C=X_0\times_{\mathbf P^1}C$, $M_C=A\times C$, and let $g:X_C\to X_0$
and $k_C:X_C\hookrightarrow M_C$ be the maps. On algebraic middle classes,
\[
 \Alg^2(X_C)\cap\ker k_{C*}
 =g^*(\Alg^2(X_0)\cap\ker k_{0*})\ \mathbin\perp\
   (\Alg^2(X_C)\cap\ker k_{C*}\cap\ker g_*).
 \tag{C44.C01.1}
\]
The first summand is canonically the old axis space
$A_B=\ker(\kappa_*:\Alg^1(B)\to H^6(A)(1))$ and has form
\[
          (w,w')\longmapsto-b\int_Bww'.                 \tag{C44.C01.2}
\]
It is positive definite, by the surface Hodge index theorem, without 032.
Consequently the remaining H9-T18 sign is exactly the sign on the second
summand in (C44.C01.1). Statements that generic covers add no algebraic
middle classes refer to this new part (or the full trace kernel), and do
not say that the entire H9-T18 variable space vanishes.
\end{proposition}
\begin{proof}
The finite flat map $g$ has $g_*g^*=b$. Projection formula makes
$P=b^{-1}g^*g_*$ a self-adjoint algebraic idempotent. The Cartesian
incidence square gives
$k_{C*}g^*=(1_A\times\pi)^*k_{0*}$, and proper functoriality gives
$k_{0*}g_*=(1_A\times\pi)_*k_{C*}$. Thus $P$ preserves
$\ker k_{C*}$, and its image there is exactly the first summand of
(C44.C01.1); its kernel is the displayed second summand. Both are
algebraic subspaces, since $P$ and $1-P$ are algebraic. Self-adjointness
proves orthogonality.

For the old space use the codimension-two blowup decomposition from
H9-T7. A middle class has unique form $q^*z+e_*p^*w$ with
$z\in H^4(A)$ and $w\in H^2(B)$. With
$t=[\mathrm{point}]\in H^2(\mathbf P^1)$, incidence pushforward is
\[
 k_{0*}(q^*z+e_*p^*w)=(hz+\kappa_*w)+zt.
\]
The exceptional divisor is $B\times\mathbf P^1$ as a variety over
$A\times\mathbf P^1$; this proves the second term's formula, while
the first follows from $[X_0]=h+t$. Kunneth independence forces $z=0$
and $\kappa_*w=0$ in the kernel. This is the asserted axis identification.
H9-T7 equation (H9.24) gives its old form $-\int_Bww'$. Pullback under
$g$ multiplies intersections by $b$, proving (C44.C01.2). Its strict
positivity is precisely the surface input used in H9-T8. No sign input
on $X_C$ has been used.
\end{proof}

\paragraph{Archive and predecessor implications.}
This sharpens the wording of the cover gate for cote 44, Th.~5,
scan 004 and JAC1964 QH9-cover-primitivep. It uses H9-T7, H9-T8,
H9-T18 and finite-flat push-pull; it does not re-prove their height
identification. Cote 16, scan 050, language is algebraic Gysin and
restriction with their correct Lefschetz weights. In Chow motives,
$b^{-1}g^*g_*$ is the old-part projector; the unresolved piece is its
complement inside the primitive middle projector.

\subsection{C44-C02: every nonextremal fixed splitting stratum of triple covers}
\begin{theorem}[C44-C02]\label{thm:maroni}
Fix the abelian pencil of H9-T23 and integers
\[
                  1\le e\le f<2e,
                  \qquad g=e+f-2\ge2.
\]
In every characteristic, let $U_{e,f}$ be the parameter space of smooth
connected triple covers of the fixed target $\mathbf P^1$, unramified
above the pencil's critical set, with framed Tschirnhausen bundle
$\mathcal O(e)\oplus\mathcal O(f)$. It is a nonempty irreducible open
subset of the space of twisted binary cubics. At its geometric generic
point, every algebraic class in the new middle trace kernel of the
pullback pencil is zero. The exact norm-height formula and radical are
\[
 h_L(z,w)=\tfrac13h_K(\operatorname{cor}z,\operatorname{cor}w),
 \qquad \operatorname{rad}h_L=
       \operatorname{cor}^{-1}(\operatorname{rad}h_K).       \tag{C44.C02.1}
\]
For abelian-fourfold pencils the full homological height is therefore
semipositive, conditionally on 032/Milne as in H9-T8. The vanishing
assertion and norm identity themselves do not use 032.

For $d=f-e>0$ the unframed stratum has codimension $d-1$ in the full
fixed-target smooth triple-cover stack. Thus this treats in particular
the first unbalanced stratum for even $g\ge6$ (a divisor), and for odd
$g\ge9$ (codimension two). It does not assert anything for every closed
point of these strata or for the extremal splitting $f=2e$.
\end{theorem}
\begin{proof}
For $E=\mathcal O(e)\oplus\mathcal O(f)$, Wood's cubic-algebra
equivalence, in the exact form already used in H9-T24 equations
(H9.136)--(H9.139), gives the affine coefficient space
\[
 P_E=H^0\bigl(\mathcal O(2e-f)\oplus\mathcal O(e)
                \oplus\mathcal O(f)\oplus\mathcal O(2f-e)\bigr).
\]
All degrees are nonnegative, so $\dim P_E=2g+8$. The smooth,
geometrically connected, and specified unramified conditions define an
open $U_E$. Each point has the required quotient bundle by the cubic
equivalence. It remains to show nonemptiness and the monodromy property.

Set
\[
             a=2e-f-1\ge0,\qquad c=f-e\ge0.
\]
Then $2a+3c=g$, $a+c+1=e$, and $a+2c+1=f$. Apply H9-T23 with
these particular parameters (its proof allows arbitrary nonnegative
$a,c$, not only the balanced choices used later in H9-T24).
In characteristic different from three, the Kummer normalization has
the relative eigensheaves
$\mathcal O\oplus\mathcal L^{-1}\oplus\mathcal L^{-2}(D_2)$,
where $\deg\mathcal L=e$ and $\deg D_2=a+1$. Its quotient dual is
$\mathcal O(e)\oplus\mathcal O(2e-a-1)=E$.
All exponent-two points and the line bundle are fixed over the moving
branch parameter, so this is an isomorphism of relative bundles.
In characteristic three, there are $a+1$ simple poles and $c$ double
poles. The exact normalized basis calculation (H9.144) gives quotient
dual degrees $a+c+1=e$ and $a+2c+1=f$, again with constant pole
divisors. After choosing a framing, in each case this constructs an
actual morphism from the H9-T23 parameter curve into $U_E$.

The universal pullback pencil is smooth proper, and
$1-\frac13Q^*Q_*$ defines a lisse middle-cohomology direct summand,
exactly as in H9-T24 equations (H9.146)--(H9.147). Its restriction to
the constructed curve has no nonzero finite-orbit vectors by H9-T23's
middle-convolution or Fourier calculation. A finite-orbit vector for
the fundamental group of $U_E$ would remain such after restriction to
this curve. Hence none exists on $U_E$. Every algebraic class at the
geometric generic point is defined over a finite extension of the
parameter function field and consequently has a finite monodromy
orbit. Its new projection must vanish. The actual-cycle push-pull and
admissibility argument in H9-T24 then gives (C44.C02.1) and the radical.
This uses precisely the old result, with the newly specified parameter
stratum, rather than a generic conclusion specialized to a closed point.

For codimension, $\dim\operatorname{Aut}(E)=d+3$ if $d>0$ and is
$4$ if $d=0$: compute $H^0(\End E)$ from its four line summands.
Thus the $d>0$ stratum has dimension $2g+5-d$, while the balanced
stratum has dimension $2g+4$. Smooth separable covers have finite
automorphism groups, so the quotient-stack dimension calculation is
valid also in characteristics two and three. The difference is $d-1$.
The first nonbalanced value is $d=2$ for even $g$ and $d=3$ for odd
$g$; $f<2e$ gives exactly the genus bounds stated above.
\end{proof}

\paragraph{Sources and scope.}
Wood, \emph{Parametrizing quartic algebras over an arbitrary base},
Theorem 2.1 and equation (1), printed pp.~1073--1074; JAC1964 H9-T23,
(H9.127)--(H9.132), and H9-T24, (H9.136)--(H9.147). These are the
precise predecessor inputs. The new conclusion concerns geometric
generic points of special Maroni strata, over algebraic closures of
their parameter fields. It is not a proof for all special covers over
$\overline{\F}_p$. Archive target: cote 44, scan 004, Th.~5 and its
fibration margin; predecessor row QH9-cover-primitivep is narrowed in
this specified stratum scope, not closed.

% END work/covers/contribution.tex

% BEGIN work/arithmetic/contribution.tex
% C44_POS, arithmetic continuation, /root/fermat_aj.
% GPT-6 Astra Pro; effort setting not exposed. 2026-11-22 (Europe/Berlin).
% New proofs below are subject to the recorded whole-proof check.

\section{The Fermat regulator: an integral projector and the exact
unramified obstruction}

The notation in this section is that of JAC1964 JAC-P13, JAC-P17,
AR-16a and AR-23.  In particular, $J/\mathbf Q$ is the Jacobian of
$x^4+y^4=z^4$, $N=h^3(J)^G$ with $G=\mu_4^2$, and
$\eta\in V:=\operatorname{CH}^2(N)_{\mathbf Q}$ is the specified
Fourier--Ceresa class.  The nonzero vector
$\alpha_5(\eta)$ spans
$S_5=H^1_f(\mathbf Q,H^3(N_{\overline{\mathbf Q}},\mathbf Q_5)(2))$
by the carried input AR-23, using AR-19 and H9-T15.  We do not reprove
that Selmer calculation, nor its nonzero cycle.  None of the deductions
in this section uses 032 or Milne's report.

The question $\ker\alpha_5=0$ remains unresolved.  The findings here
locate its integral obstruction more precisely.  They do not identify
a Selmer dimension with a Chow dimension.

\subsection{An integral model at the prime 5}

\begin{proposition}[C44-F01]\label{C44-F01}
Let $R=\mathbf Z_{(5)}$ and $S=\operatorname{Spec}\mathbf Z[1/10]$.
Let $\mathcal J/S$ be the Jacobian of the smooth Fermat curve
$\mathcal C:x^4+y^4=z^4$.  There is an exact idempotent, with
coefficients in $\mathbf Z[1/2]\subset R$, given by
\[
 e=\frac{1-[-1]^*}{2}\,\frac1{16}\sum_{g\in G}g^*.
 \tag{F01.1}
\]
The sum is formed after $S'=\operatorname{Spec}\mathbf Z[i,1/10]$
and descended to $S$.  Its rational generic-fibre motive is exactly
$N$.  Put
\[
 A=e\operatorname{CH}^2(J;R).
\]
Then restriction is an isomorphism
\[
 e\operatorname{CH}^2(\mathcal J;R)
       \xrightarrow{\ \sim\ } A,
 \qquad A\otimes_R\mathbf Q=V.                 \tag{F01.2}
\]
Moreover $e$ acts on geometric $5$-adic cohomology only in degree
three, where its image is a free rank-two direct summand.  Thus
\[
 T=eH^3_{\mathrm{et}}(J_{\overline{\mathbf Q}},\mathbf Z_5)(2)
 \tag{F01.3}
\]
is an actual $\mathbf Z_5$ lattice for the representation in the task,
and $eH^b_{\mathrm{et}}(J_{\overline{\mathbf Q}},\mathbf Z/5^n(2))=0$
for $b\ne3$ and every $n\ge1$.
\end{proposition}

\begin{proof}
The Fermat equation is smooth over $\mathbf Z[1/2]$: a singular
projective geometric point would have $x=y=z=0$, since all three
partial derivatives are units times the corresponding cubes.
The section $(1:0:1)$ gives its relative Jacobian, an abelian scheme.
The group action and inversion extend over $S$.  The two averaging
operators in (F01.1) commute and are idempotent.  They use only
powers of two as denominators.  Galois conjugation on $S'/S$ permutes
the sixteen group elements.  Restriction and norm, whose composite
is multiplication by two, descend the idempotent with the asserted
coefficients.  This descent is valid for Chow correspondences and
also for their actions on the cohomology theories used below.

The six characters of $G$ on $H^1(J)$ are
\[
 (1,1),(1,2),(2,1),(3,3),(3,2),(2,3)\quad\text{in }(\mathbf Z/4)^2;
 \tag{F01.4}
\]
these are the carried calculation in JAC-P17.  In their exterior
powers, an odd-degree invariant occurs only in degree three.  The
two invariant triples are the first three and the last three
characters.  Their types are $(3,0)$ and $(0,3)$.  The exact rational
identification with the projector of $N$ is JAC-P13/JAC-P17; in
particular (F01.1) is an integral-at-5 representative of the same
rational summand, not a new motive inferred solely from its
realization.  More explicitly, averaging is zero on $h^1(J)$
because it is induced by the zero norm endomorphism
$\sum_{g\in G}g$.  It is zero on $h^5(J)$ by the duality
identification through the $G$-invariant principal polarization.
Inversion removes the even Deninger--Murre pieces.  Thus the
rational idempotent on the full motive equals $\pi_3p_G$.
The geometric cohomology of an abelian variety is the
exterior algebra of its free degree-one cohomology.  Because the
projector has denominators prime to five, its integral image is a
direct summand and has no torsion.  The character calculation
therefore proves the assertions with both $\mathbf Z_5$ and
$\mathbf Z/5^n$ coefficients.

We next verify that $e$ annihilates the vertical terms in Chow
localization.  On any fibre of $\mathcal J$, the operator
$(1-[-1]^*)/2$ sends divisor classes into
$\operatorname{Pic}^0\otimes R$.  The endomorphism
$\sum_{g\in G}g$ of $\mathcal J_{S'}$ is zero.  Indeed its generic
fibre has zero tangent map by (F01.4), so is zero, and equality of
endomorphisms of an abelian scheme is detected on its dense generic
fibre.  Its dual is also zero, which says that
$\sum g^*$ annihilates $\operatorname{Pic}^0$.  Consequently
\[
 e\operatorname{CH}^1(\mathcal J_s;R)=0
 \quad\text{for every closed }s\in S.            \tag{F01.5}
\]
When a fibre does not contain $i$, check this after the extension
of degree at most two and descend using restriction--norm.

Chow localization gives
\[
 \bigoplus_{s\in S^{(1)}}\operatorname{CH}^1(\mathcal J_s;R)
 \longrightarrow\operatorname{CH}^2(\mathcal J;R)
 \longrightarrow\operatorname{CH}^2(J;R)\longrightarrow0.
 \tag{F01.6}
\]
The action of every automorphism in (F01.1) commutes with these
maps.  Applying the exact direct-summand functor $e$ and (F01.5)
proves (F01.2).  Notice that this argument is integral at five;
it does not invoke rational vanishing of the whole Chow group of
a special fibre.
\end{proof}

\paragraph{Locators and scope.}
The localization sequence is Fulton, \emph{Intersection theory},
Proposition~1.8 and its arithmetic-base version in \S20.2.
The Picard-scheme construction is FGA, Expos\'e~232,
Theorem~3.1 (printed pp.~148--149), applied to the smooth
projective relative curve; the canonical abelian component is
FGA, Expos\'e~236, Theorem~3.5 (printed p.~239), with the smooth
Picard schemes of the curve fibres.  The latter fibre smoothness
also follows from the vanishing of their $H^2(\mathcal O)$.
The relative-motive formalism over schemes smooth over a Dedekind
domain is recalled in the supplied Laga--Shnidman source,
\emph{Ceresa cycles of bielliptic Picard curves}, \S4.1--4.5;
the extension of Deninger--Murre to that base is explained at the
start of \S4.2, citing K\"unnemann, \S1, Remark~1.1.
The generic projector, the character list, and its rational field
of definition are JAC1964 JAC-P13 and JAC-P17.  The new assertion
here is the integral-at-5 localization isomorphism for that
projector.  The finite character check is reproduced by
\texttt{work/arithmetic/check\_fermat\_projector.py}; the script
does not prove any statement about rational equivalence.

\begin{lemma}[C44-F02]\label{C44-F02}
For the lattice $T$ in (F01.3),
\[
 H^0(\mathbf Q,T/5^nT)=0\quad(n\ge1),
 \qquad H^0(\mathbf Q,T\otimes\mathbf Q_5/\mathbf Z_5)=0.
 \tag{F02.1}
\]
Consequently
\[
 eH^3_{\mathrm{et}}(J,\mathbf Z/5^n(2))=0,
 \qquad
 eH^4_{\mathrm{et}}(J,\mathbf Z/5^n(2))
       \simeq H^1(\mathbf Q,T/5^nT).             \tag{F02.2}
\]
\end{lemma}

\begin{proof}
We make AR-23's Frobenius-at-11 check independent of a choice of
integral lattice.  By JAC-P17, the weight-three character is
$\psi_+^3$, where $\psi_+$ is the character of
$E_+:y^2=x^3+x$.  Direct counting gives
$\#E_+(\mathbf F_{11})=12$, so the degree-one Frobenius
polynomial is $X^2+11$.  Cubing its two roots gives
$X^2+11^3$ on $N$.  After the twist $(2)$, the polynomial is
\[
                 X^2+\frac1{11}.
\]
It follows that
\[
       \det(1-\operatorname{Fr}_{11}\mid T)
              =\frac{12}{11}\in\mathbf Z_5^\times.
 \tag{F02.3}
\]
For any stable $\mathbf Z_5$ lattice, this determinant condition
makes $1-\operatorname{Fr}_{11}$ invertible.  Its reductions are
invertible modulo every $5^n$, and it is invertible on the torsion
quotient as well.  Invariance under the entire Galois group implies
invariance under this Frobenius, proving (F02.1).  The conclusion
does not change if arithmetic rather than geometric Frobenius is
used, since both Frobenius and its inverse are integral units.

Apply the Hochschild--Serre spectral sequence with finite
coefficients and the integral projector $e$.  By C44-F01 the
projected geometric cohomology occurs only in row $b=3$.
The spectral sequence therefore degenerates on that row and
identifies its total degrees three and four with the two Galois
groups in (F02.1)--(F02.2).
\end{proof}

\paragraph{Check scope.}
The numerical part of (F02.3), including the point count and its
reduction $12/11\equiv2\pmod5$, is an exact integer/rational
computation in the accompanying script.  This is a lattice-uniform
form of the Frobenius test used in AR-23, not a new proof of
AR-23's nonzero Abel--Jacobi class.

\subsection{The precise obstruction after passing to finite coefficients}

Write
\[
 B=eH^4_{\mathrm{et}}(J,\mathbf Z_{(5)}(2)),
 \qquad
 C=eH^3_{\mathrm{nr}}(J,\mathbf Q_5/\mathbf Z_5(2)).
 \tag{F03.1}
\]
The cohomology defining $B$ is \emph{\'etale motivic cohomology},
not continuous $5$-adic cohomology.  The distinction is essential.
The group $C$ is unramified cohomology over the ground field
$\mathbf Q$, with residues at the codimension-one points of $J$.
The operations defining $e$ act on it by restriction, norm and
the automorphisms of $J_{\mathbf Q(i)}$.

\begin{theorem}[C44-F03]\label{C44-F03}
There is a canonical short exact sequence
\[
                 0\longrightarrow A\longrightarrow B
                   \longrightarrow C\longrightarrow0,
 \tag{F03.2}
\]
and both $A$ and $B$ have no nonzero $5$-primary torsion.
For every $n\ge1$, the kernel of the finite-coefficient cycle map
\[
 c_n:A/5^nA\longrightarrow H^1(\mathbf Q,T/5^nT)
\]
is canonically $C[5^n]$.  In particular, for
$\widehat A=\varprojlim_n A/5^nA$, the completed cycle map satisfies
\[
 \ker\!\left(\widehat c:\widehat A\longrightarrow
                    H^1_{\mathrm{cont}}(\mathbf Q,T)\right)
          = T_5C:=\varprojlim_n C[5^n],            \tag{F03.3}
\]
where the transition on $C[5^n]$ is multiplication by five.

Let
\[
 \rho:A\otimes_R\mathbf Q_5\longrightarrow
                  \widehat A\otimes_{\mathbf Z_5}\mathbf Q_5,
 \qquad K_{\mathrm{comp}}=\ker\rho.
\]
The full scalar-extended Abel--Jacobi kernel in the question has
the exact description
\[
 0\longrightarrow K_{\mathrm{comp}}
 \longrightarrow\ker\alpha_5
 \xrightarrow{\rho}
 \operatorname{im}\rho\cap
             (T_5C\otimes_{\mathbf Z_5}\mathbf Q_5)
 \longrightarrow0.                              \tag{F03.4}
\]
The intersection is taken inside
$\widehat A\otimes\mathbf Q_5$.  Neither group on the two ends
of (F03.4) has been proved to vanish.
\end{theorem}

\begin{proof}
For a smooth variety over a field, the codimension-two
comparison sequence is
\[
 0\longrightarrow\operatorname{CH}^2(J)
 \longrightarrow H^4_{\mathrm{et}}(J,\mathbf Z(2))
 \longrightarrow H^3_{\mathrm{nr}}(J,\mathbf Q/\mathbf Z(2))
 \longrightarrow0.
\]
The exact locator is Kahn, \emph{On the injectivity and
non-injectivity of the $l$-adic cycle class maps}, (1.5),
author p.~3, citing his \emph{Classes de cycles motiviques
\'etales}, Proposition~2.9.  Localizing at five and applying
the idempotent $e$ gives (F03.2).

The coefficient triangle for $\mathbf Z_{(5)}(2)$ has
$\mathbf Z/5^n(2)=\mu_{5^n}^{\otimes2}$ as its finite-coefficient
term.  Its long exact sequence and (F02.2) give
$B[5^n]=0$ and an injection
\[
 B/5^nB\hookrightarrow
       eH^4_{\mathrm{et}}(J,\mathbf Z/5^n(2)).     \tag{F03.5}
\]
The relevant coefficient comparison is the Merkurjev--Suslin
case of the norm-residue theorem; it is recorded in Kahn,
loc.~cit., (3.2) and (4.1), author pp.~7 and~10.
Since $A\hookrightarrow B$, also $A[5^n]=0$.

Apply multiplication by $5^n$ to (F03.2).  The snake lemma gives
\[
 0\longrightarrow C[5^n]\xrightarrow{\delta_n}A/5^nA
 \longrightarrow B/5^nB.
 \tag{F03.6}
\]
Explicitly, for $c\in C[5^n]$, choose $b\in B$ above $c$;
then $5^n b\in A$, and
$\delta_n(c)=5^n b\bmod5^nA$.  This is independent of the
choice of lift.  The absence of $5^n$-torsion in $B$ proves
its injectivity.  Composing with (F03.5) and the
Hochschild--Serre identification gives precisely $c_n$.
Thus (F03.6) proves the finite-coefficient assertion.

If $c\in C[5^{n+1}]$ and $b$ is a lift, then the lift $5b$
of $5c$ shows that reduction of $\delta_{n+1}(c)$ modulo $5^n$
equals $\delta_n(5c)$.  The transition is therefore multiplication
by five.  Taking inverse limits in (F03.6) is exact at the first
two terms, so the kernel of the completed finite-coefficient
cycle map is $T_5C$.  This uses left exactness only; no
unverified surjectivity of transitions is needed.

Continuous cohomology has the comparison exact sequence
\[
 0\longrightarrow\varprojlim{}^1_n
       H^0(\mathbf Q,T/5^nT)
 \longrightarrow H^1_{\mathrm{cont}}(\mathbf Q,T)
 \longrightarrow\varprojlim_n H^1(\mathbf Q,T/5^nT)
 \longrightarrow0.
\]
Its left term is zero by C44-F02.  Equivalently, in total
degree four on $J$ the potential $\varprojlim{}^1$ term is the
limit of the zero groups in (F02.2).  Thus the completed map
just constructed is the continuous cycle map, proving (F03.3).
This is why an inverse limit is legitimate in this particular
summand although it is not interchangeable with continuous
cohomology for an arbitrary variety over a number field.

After tensoring by $\mathbf Q_5$, the cycle map from $A$
factors as $\widehat c\circ\rho$.  Its rational projected
Hochschild--Serre component is the Abel--Jacobi map
$\alpha_5$ from the task.  The inclusion of its image in
$H^1_f$ is the carried AR-16a formalism and does not change
the kernel.  Flatness of $\mathbf Q_5$ over $\mathbf Z_5$
identifies the kernel of $\widehat c\otimes\mathbf Q_5$ with
$T_5C\otimes\mathbf Q_5$.  Taking the inverse image under
$\rho$ and then its image yields (F03.4).
\end{proof}

\paragraph{What this attack removes, and what it leaves.}
C44-F01 removes the ambiguity of a varying arithmetic spread for
the Chow source.  C44-F02 removes both the global degree-three
term and the continuous-cohomology $\varprojlim{}^1$ ambiguity
from the finite-coefficient calculation.  The remaining group
in (F03.3) is the actual Tate module of one specified unramified
group.  It is not the Selmer group.  A Selmer bound does not bound
this Tate module or prove the tensor-to-completion injectivity in
(F03.4).  Specialization of rational odd-motive Chow classes,
already ruled out as a detector by AR-20, supplies neither
missing assertion.

\subsection{A bounded-denominator formulation on rational planes}

For a finite free $R$-submodule $M\subset A$, let
$W_M=M\otimes_R\mathbf Q\subset V$.  Equation (F03.2) and
the absence of $5$-torsion identify $B$ with a subgroup of $V$:
if $b\in B$, choose $5^n b=a\in A$ and send $b$ to $a/5^n$.
This map is well defined and injective, since $C$ is $5$-primary
torsion.  This convention makes $B\cap W_M$ unambiguous.

\begin{proposition}[C44-F04]\label{C44-F04}
The following are equivalent for every finite free $M\subset A$:
\begin{enumerate}
\item $\alpha_5$ is injective on $M\otimes_R\mathbf Q_5$;
\item $T_5(B/M)=0$;
\item $(B\cap W_M)/M$ is finite;
\item there is an integer $b_M\ge0$ such that
      $5^{b_M}(B\cap W_M)\subset M$.
\end{enumerate}
More precisely,
\[
 \ker(\alpha_5|_{M\otimes\mathbf Q_5})
       \simeq T_5(B/M)\otimes_{\mathbf Z_5}\mathbf Q_5.
 \tag{F04.1}
\]
Choose a nonzero multiple $a_\eta\in A$ of the carried class
$\eta$.  It is enough to verify these equivalent conditions for
the modules
\[
                  M_z=R a_\eta+Rz\quad(z\in A).
 \tag{F04.2}
\]
The condition holds when $z$ lies in $\mathbf Q\eta$, by AR-23.
For arbitrary $z$ the bounded-denominator condition in (4) is
unproved.
\end{proposition}

\begin{proof}
Apply the argument of (F03.6) to
$0\to M\to B\to B/M\to0$.  Both $M$ and $B$ have no
$5$-torsion.  It gives
\[
 0\longrightarrow T_5(B/M)\longrightarrow\widehat M
                         \longrightarrow\widehat B.
\]
Since $M$ is finite free over $R$, its completion is
$M\otimes_R\mathbf Z_5$.  The injection (F03.5) implies
that the kernel of the map to $\widehat B$ is exactly the
kernel of the continuous cycle map.  Rationalizing proves
(F04.1).

All $5$-power-torsion elements of $B/M$ lie in
\[
 D_M=(B\cap W_M)/M\ \subset\ W_M/M
          \simeq(\mathbf Q_5/\mathbf Z_5)^{\operatorname{rank}M}.
\]
Conversely every element of $D_M$ is $5$-power torsion, since
$R$ already inverts the other primes.  Thus $T_5(B/M)=T_5D_M$.
A subgroup of a finite direct sum of $\mathbf Q_5/\mathbf Z_5$
is a direct sum of finitely many copies of
$\mathbf Q_5/\mathbf Z_5$ and a finite group.  One proof is to
take its Pontryagin dual: it is a quotient of the finite free
$\mathbf Z_5$ module, hence is finitely generated; the structure
theorem over the discrete valuation ring $\mathbf Z_5$ and
duality give the assertion.  It follows that $T_5D_M$ is finite
free, and vanishes exactly when $D_M$ is finite.  Finiteness of
this subgroup is equivalent to a bound on its exponent, giving
the equivalence with (4) and also with (1), by (F04.1).

Finally AR-23 makes $\alpha_5(a_\eta)\ne0$ and gives a
one-dimensional target.  If all the two-generator modules
(F04.2) have injective scalar-extended map, none can have rank
two.  Hence every $z\in A$ belongs to $\mathbf Q\eta$, so
$V=\mathbf Q\eta$ and the full map is injective.  The converse
is immediate.  On the line $\mathbf Q\eta$, AR-23 already gives
injectivity, so the equivalent bounded-denominator conclusion
is proved on that line without any hypothesis about the rest
of the Chow group.
\end{proof}

\begin{proposition}[C44-F05: an algebraic model of the missing step]
\label{C44-F05}
Even finite generation of the integral Chow-like source,
one-dimensional surjective regulator image, a nonzero distinguished
class, and a positive definite real form do not formally imply
the scalar-extended injectivity in C44-F04.  There is an explicit
model satisfying all those properties in which the completed
kernel is one-dimensional and arises from a nonzero unramified-like
Tate module.
\end{proposition}

\begin{proof}
Choose the root $a\in\mathbf Z_5$ of $a^2=6$ with $a\equiv1\pmod5$.
Hensel's lemma gives this root, and $a\notin\mathbf Q$ because
six is not a rational square.  Put
\[
 A_0=R^2,\qquad
 B_0=\{(r,s)\in\mathbf Q^2:r+as\in\mathbf Z_5\},
 \qquad f(r,s)=r+as.
\]
Then $A_0\subset B_0$, both are $5$-torsion-free, and
\[
  B_0/A_0\simeq\mathbf Q_5/\mathbf Z_5.
 \tag{F05.1}
\]
To see the last statement, take the kernel of the surjection
$(\mathbf Q_5/\mathbf Z_5)^2\to\mathbf Q_5/\mathbf Z_5$,
$(r,s)\mapsto r+as$, whose restriction on the first coordinate
is the identity.  The inverse image of this kernel in
$\mathbf Q^2$ is exactly $B_0$.

The scalar-extended regulator is
$\mathbf Q_5^2\to\mathbf Q_5$, $(r,s)\mapsto r+as$.
It is surjective, sends $\eta_0=(1,0)$ to one, and has kernel
$\mathbf Q_5(-a,1)$.  It is even injective on the underlying
rational vector space $\mathbf Q^2$.  The form
$h((r,s),(r',s'))=rr'+ss'$ is positive definite and
$h(\eta_0,\eta_0)=1$.  Neither fact removes the scalar-extended
kernel.

For a completely explicit division tower, choose integers $a_n$
with $a_n\equiv a\pmod{5^n}$ and
$a_{n+1}\equiv a_n\pmod{5^n}$, and set
\[
 b_n=5^{-n}(-a_n,1)\in B_0.
\]
Then $b_n\bmod A_0$ has order $5^n$ and
$5b_{n+1}-b_n\in A_0$.  The resulting nonzero element of
$T_5(B_0/A_0)$ maps to $(-a,1)\in\widehat A_0$.
Thus the precise obstruction in C44-F03 is present in the
model, despite all the other listed data.  This is a model of
an invalid inference, not a geometric counterexample to
Abel--Jacobi injectivity or to a height conjecture.
\end{proof}

\subsection{What finite coefficients see of the fourfold component}

\begin{proposition}[C44-F06]\label{C44-F06}
Let $F/\mathbf Q$ be finite and let $D/F$ be a smooth projective
geometrically connected curve with a point $d_0\in D(F)$.
On $X=D\times J_F$, put
\[
 q=\pi_1^D\otimes e,\qquad
 A_D=q\operatorname{CH}^2(X;R),\qquad
 B_D=qH^4_{\mathrm{et}}(X,R(2)),
 \quad R=\mathbf Z_{(5)}.
\]
Then $B_D$ is uniquely divisible by five and hence is a
$\mathbf Q$-vector space.  The motivic comparison identifies it
canonically with the rational Chow component:
\[
       B_D\simeq A_D\otimes_R\mathbf Q.
                                                        \tag{F06.1}
\]
In particular $A_D$ has no $5$-primary torsion, and the
finite-coefficient cycle map on $A_D/5^nA_D$ is zero for every $n$.
Writing
\[
 C_D=qH^3_{\mathrm{nr}}(X,\mathbf Q_5/\mathbf Z_5(2)),
\]
one has
\[
 C_D=B_D/A_D,\qquad
 C_D[5^n]\simeq A_D/5^nA_D,
 \qquad T_5C_D\simeq\widehat A_D.             \tag{F06.2}
\]
Thus finite generation of $B_D$ as an $R$-module would force
the entire rational component in JAC-P18 to vanish.  Such finite
generation is not established here; finite generation of $A_D$
alone would not have that consequence.
\end{proposition}

\begin{proof}
The point projector $\pi_1^D$ is integral.  By C44-F01 and
the K\"unneth formula, $q$ selects geometric cohomology only in
degree four, where the integral lattice after twist is
\[
 T_D=H^1_{\mathrm{et}}(D_{\overline F},\mathbf Z_5)
                  \otimes_{\mathbf Z_5}T.
\]
It is free, of rank $4g(D)$.  At any place of $F$ over five,
its rational realization has Hodge--Tate weights
$-2,-1,1,2$ (each with multiplicity $g(D)$, up to the
simultaneous opposite weight convention).  Indeed the two
factors before twist have weights $\{0,1\}$ and $\{0,3\}$.
None is zero.  The Hodge--Tate comparison for smooth proper
varieties therefore implies
\[
                 H^0(F,T_D)=0.                 \tag{F06.3}
\]
Here one may use Scholze, \emph{$p$-adic Hodge theory for
rigid-analytic varieties}, Corollary~1.8, author p.~4, applied
to the analytification, together with functoriality for the
algebraic projectors.  This use needs neither good reduction
of $D$ at five nor a Hodge or Tate conjecture.

Finite-coefficient Hochschild--Serre now gives
\[
 qH^3_{\mathrm{et}}(X,\mathbf Z/5^n(2))=0,
 \qquad
 qH^4_{\mathrm{et}}(X,\mathbf Z/5^n(2))
       =H^0(F,T_D/5^nT_D).
                                                        \tag{F06.4}
\]
As in C44-F03, the coefficient triangle implies that
$B_D[5^n]=0$ and injects $B_D/5^nB_D$ into the second group
in (F06.4).  Taking inverse limits gives
\[
 \widehat B_D\hookrightarrow
        \varprojlim_n H^0(F,T_D/5^nT_D)
        =H^0(F,T_D)=0.
\]
The maps $B_D/5^{n+1}B_D\to B_D/5^nB_D$ are surjective,
so each projection from their inverse limit is surjective.
Consequently $B_D/5^nB_D=0$ for every $n$.  Together with
$B_D[5^n]=0$ this proves unique divisibility by five.
The other primes are already inverted in $R$.

The comparison sequence (F03.2) is valid in codimension two
on $X$ as well, and is compatible with $q$.  Its quotient
is $5$-primary torsion.  Rationalizing its injection
$A_D\hookrightarrow B_D$, and using that $B_D$ is already
a rational vector space, proves (F06.1).  The finite-coefficient
cycle map factors through $B_D/5^nB_D=0$.  Applying the snake
lemma to multiplication by $5^n$ gives (F06.2), including the
transition by multiplication by five; its inverse limit is
the final displayed isomorphism.

A nonzero rational vector space is not finitely generated as
a $\mathbf Z_{(5)}$-module.  Thus the stated finite generation
condition forces $B_D=0$.  JAC-P18 has already proved that
every algebraic class in this projected component is
homologically zero, so this is the vanishing that its
componentwise Murre~(D) would require.  The elementary inclusion
$R\subset\mathbf Q$ shows why finite generation of the smaller
group $A_D$ cannot substitute for finite generation of $B_D$.
\end{proof}

\paragraph{Scope of the additional attack.}
C44-F06 does not prove Murre~(D).  It identifies the relevant
\'etale motivic group as entirely divisible, with all of its
completed Chow source accounted for by unramified cohomology.
This makes precise why torsion-freeness or finite-coefficient
cycle computations alone cannot kill the JAC-P18 obstruction.
It concerns cote~44, Th.~1--3 (scan~002) and the boundary of
Th.~5 (scan~004); it does not change the known polarization
on the incidence higher Picard factor.  In the frameworks of
SGA and Chow motives, it is the tensor summand
$h^1(D)\otimes N$ of weight four with no $(2,2)$ part,
followed by the same codimension-two \`etale-motivic comparison
as C44-F03.  Its vanishing is a rational-equivalence issue,
not the semisimplicity or signed $I$ assertion of cote~16,
scan~049.  No 032/Milne input is used.

\subsection{The remaining finiteness question on one arithmetic model}

\begin{proposition}[C44-F07]\label{C44-F07}
For the abelian scheme $\mathcal J/S$ of C44-F01, restriction
also gives an isomorphism on the projected \emph{\'etale
motivic} groups:
\[
 eH^4_{\mathrm{et}}(\mathcal J,R(2))
                 \xrightarrow{\ \sim\ } B.
                                                        \tag{F07.1}
\]
The completion $\widehat B$ is a free $\mathbf Z_5$-module
of rank exactly one.  Its rational realization is the Selmer
line already realized by $\eta$.

The following two conditions are equivalent:
\begin{enumerate}
\item the scalar-extended Fermat Abel--Jacobi map is injective;
\item $eH^4_{\mathrm{et}}(\mathcal J,R(2))$ is finitely
      generated over $R$.
\end{enumerate}
Condition (2) is not proved here.  In particular the proved
finite generation of the completion is not being substituted
for it.
\end{proposition}

\begin{proof}
Purity and localization for \'etale motivic cohomology
give the exact segment
\[
 \bigoplus_{s\in S^{(1)}}
 H^2_{\mathrm{et}}(\mathcal J_s,R(1))
 \longrightarrow H^4_{\mathrm{et}}(\mathcal J,R(2))
 \longrightarrow H^4_{\mathrm{et}}(J,R(2))
 \longrightarrow\bigoplus_{s\in S^{(1)}}
 H^3_{\mathrm{et}}(\mathcal J_s,R(1)).
                                                        \tag{F07.2}
\]
An exact source for the needed purity is Geisser,
\emph{Motivic cohomology over Dedekind rings},
Theorem~1.2(1), published p.~774: after inverting the residue
characteristic no truncation is needed.  All residue
characteristics of $S$ differ from five and are units in $R$.
The Bloch--Kato premise of that 2004 statement is supplied by
the norm-residue theorem (in this weight, Merkurjev--Suslin).
The localization of the cycle complex and passage to the
generic fibre give the displayed direct sum.  This is also
the localization used in Geisser, \emph{Hasse principles for
\'etale motivic cohomology}, \S4, diagram~(2).

The first outer group in (F07.2) is a sum of Picard groups,
because $\mathbf Z(1)_{\mathrm{et}}=\mathbf G_m[-1]$.
Its $e$-part vanishes by C44-F01.  The last is a sum of
cohomological Brauer groups localized at five.  For a fibre
$\mathcal J_s$, finite-coefficient Hochschild--Serre and the
character calculation of C44-F01 give
\[
           eH^2_{\mathrm{et}}(\mathcal J_s,\mu_{5^n})=0
           \quad(n\ge1).
\]
Indeed its projected geometric cohomology starts in degree
three.  Kummer's sequence implies
$e\operatorname{Br}(\mathcal J_s)[5^n]=0$ for every $n$.
The Brauer group is torsion, so its localization at five has
zero $e$-part.  Apply $e$ to (F07.2) to obtain (F07.1).

The lattice $T$ is a lisse $\mathbf Z_5$-sheaf on $S$.
Proper smooth base change and the same single-row calculation
give
\[
 eH^4_{\mathrm{et}}(\mathcal J,\mathbf Z/5^n(2))
          =H^1_{\mathrm{et}}(S,T/5^nT),
 \qquad eH^3_{\mathrm{et}}(\mathcal J,\mathbf Z/5^n(2))=0.
                                                        \tag{F07.3}
\]
For the second identity, Frobenius at eleven still belongs
to $\pi_1(S)$, so C44-F02 applies.  Finite-coefficient
cohomology of this arithmetic scheme is finite.  Alternatively,
in degree one this follows from the finiteness of number
field extensions of bounded degree unramified outside
$\{2,5,\infty\}$.
It follows from the coefficient comparison that each
$B/5^nB$ is finite, and that
\[
 \widehat B\hookrightarrow
           H^1_{\mathrm{cont}}(S,T).
                                                        \tag{F07.4}
\]
There is no inverse-limit correction by (F07.3).
The group on the right is a finite free $\mathbf Z_5$-module:
it is a compact inverse limit of finite groups, its quotient
modulo five injects into $H^1(S,T/5T)$, and hence compact
Nakayama gives finite generation; its $5$-torsion vanishes
because $H^0(S,T/5T)=0$.  Restriction to
$H^1_{\mathrm{cont}}(\mathbf Q,T)$ is injective by inflation.

The group $B$ is rationally $V$, by (F03.2).  Therefore the
rational image of every element of $B$ in (F07.4) belongs to
the Selmer line of AR-23.  Its intersection with the finite
free lattice $H^1_{\mathrm{cont}}(S,T)$ is a closed rank-one
lattice.  The image of $B$ is dense in $\widehat B$, whose
image in (F07.4) is closed by compactness.  Hence
$\widehat B$ lies in that rank-one lattice.  It is nonzero,
because $a_\eta\in A\subset B$ has nonzero regulator.
This proves its asserted rank, including its rational image.

If $B$ is finitely generated over $R$, torsion-freeness makes
it finite free.  Its rank equals the rank of $\widehat B$,
namely one.  As $A\subset B$ and $A\otimes_R\mathbf Q=V$,
this forces $V=\mathbf Q\eta$, and AR-23 then proves
injectivity.  Conversely, injectivity and AR-23 give
$V=\mathbf Q\eta$.  Apply C44-F04 to $M=Ra_\eta$: since
its rational span is all of $V$, the bounded-denominator
conclusion places $B$ inside $5^{-b_M}M$.  Thus $B$ is a
submodule of a finite free $R$-module and is finitely
generated.  This proves the equivalence.
\end{proof}

\paragraph{Meaning of the arithmetic finiteness gate.}
C44-F07 makes the finiteness question a statement about one
specific summand on one scheme over $\mathbf Z[1/10]$.
Neither finite generation of ordinary Chow groups nor
finite generation of continuous \(5\)-adic cohomology proves
it.  Kahn, \S4, author pp.~10--11, explicitly isolates this
distinction; his Proposition~4.1 gives counterexamples to
unqualified \`etale-motivic finiteness on arithmetic schemes.
Geisser, \emph{Hasse principles}, Conjecture~4.3, states
Lichtenbaum's finiteness expectation for regular schemes
proper over $\mathbf Z$; that conjecture is not a theorem
for the present summand, and $\mathcal J/S$ is proper over
the indicated open base, not over all of $\mathbf Z$.
No claim of direct applicability or proof of that conjecture
is made.  In cote~44 this refines the arithmetic gate of
Th.~5, scan~004; in cote~16 it specifies information about
rational equivalence that conditions $C$ and $I$ on pure
realizations do not provide.  The equivalence does not by
itself count as a strictly smaller main gate under the
progress rule of the prompt.  No 032/Milne input is used.

\subsection{The three proposed arithmetic routes: exact scope}

\paragraph{C44-FA01: Euler systems.}
The supplied Lilienfeldt--Shnidman paper,
\emph{Derivatives of Rankin--Selberg $L$-functions and heights of
generalized Heegner cycles}, arXiv:2408.04375,
Theorem~1.1/13.4, identifies a height on a Hecke-generated cycle
space modulo the real-height radical.  Its Corollary~1.4 expressly
assumes the injectivity of its scalar-extended Abel--Jacobi map.
In \S15, author pp.~71--73, equation~(99) on p.~72 is the Euler-system
conclusion that a nonzero generalized Heegner cohomology class
spans the Selmer space; the next paragraph invokes the assumed
injectivity to pass back to Chow.  The displayed implication chain
on p.~73 marks both uses of Abel--Jacobi injectivity.  That is the exact missing
step, and C44-F03--F04 describe its integral obstruction for the
Fermat summand.  No Euler-system relation at the level of
Galois cohomology bounds the denominators in (F04.2).

There are also independent obstacles to a direct substitution
in the displayed theorem: its standing assumptions in \S1.1
require an imaginary quadratic field of odd discriminant and
an unramified Hecke character, whereas the Fermat character
over $\mathbf Q(i)$ has conductor $(1+i)^4$ by JAC-P17/AR-17.
The stated range is $0<t\le k-1$, so the weight-four central
case with $k+t=2$ is not a direct instance.  These scope failures
are not the reason the existing Selmer result fails---that result
is already available as AR-19/H9-T15---but they prevent the paper
from being treated as a new Chow-generation theorem here.

The Loeffler--Zerbes constructions consulted in this run
interpolate or bound Galois-cohomology classes, not the entire
Chow group.  The exact example is Jetchev--Loeffler--Zerbes,
\emph{Heegner points in Coleman families}, Theorem~A
(published p.~125), equivalently Theorem~5.4.1, whose
specializations are \emph{Abel--Jacobi images} of generalized
Heegner cycles.  An interpolation statement about these images
does not supply the missing denominator bound.  The existing
JAC1964 audit AR-17 and its
\texttt{LATEST\_RANK\_SOURCE\_SCOPE.tex} remain in force; their
statements are not reproved here.

\paragraph{C44-FA02: Gross--Schoen heights.}
AR-24--AR-25 are carried as the exact projected-diagonal and
sixteen-intersection identities.  Zhang,
\emph{Gross--Schoen cycles and dualising sheaves},
Theorem~2.3.5, author p.~27, proves the triple-intersection
identity; Conjecture~2.4.1, author p.~28, asserts semipositivity
of the tensor-correspondence pairing.  The latter is not proved
by the former.  Even adding that conjecture would give the
semipositivity conclusion recorded in AR-25, without an
equality criterion or the injectivity in task~3.  In particular
one must not substitute a real-height calculation for
the absent faithfulness of the $5$-adic regulator.  C44-F05
shows that positivity and the stated Selmer data alone are not
a formal bridge.  The supplied Gu draft, Theorem~1.3,
Corollary~5.4 and Conjecture~1.1, gives formulas and
conditional positivity, not that missing faithfulness.  Its
statements are treated here as supplied draft inputs, not as
independently verified public results, and none is used in the
proofs of C44-F01--F05.

\paragraph{C44-FA03: the algebraically trivial radical.}
JAC-P18 isolates the possible cycle
$Z=(\pi_1^C\otimes\pi_N)_*Z$ on $C\times J$, fixed by the
degree-four product projector.  Xu--Xu,
\emph{Remarks on Murre's conjecture on Chow groups},
Corollary~3.9, author p.~9, gives (A) and (B) for the product;
the earlier definition of (D), author p.~2, is the extra
cycle-class injectivity needed to annihilate this component.
Theorem (B) leaves degree four acting on codimension two and
therefore does not remove this component.  The special
Laga--Shnidman isomorphism in their Theorem~5.2 concerns a
different invariant summand of types $(2,1),(1,2)$, identified
with $h^1(E)(-1)$.  Such a Picard-motive isomorphism cannot hold
for the Fermat $N$, whose types are $(3,0),(0,3)$.
No proof of the needed componentwise (D) was found in these
sources or in the other primary sources checked in this run.

\paragraph{Programme crosswalk and current gate.}
C44-F01--F04 concern the arithmetic extension of cote~44,
scan~004, Th.~5, and the incidence/higher-Picard boundary of
Th.~1--3 on scan~002, as made explicit in the letter on
scans~008--009.  They do not alter JAC-T4.3 or establish a new
polarization of $P^i(X)$.  The actual cycle lives in the
extreme-type motive outside that Picard component.
In cote~16, scans~049--050, the conclusion is a limitation of
the passage from $C$, $I$, and semisimplicity of cohomological
correspondences to rational Chow groups: the kernel measured
by (F03.4) is invisible to those pure-realization assertions.
The corrected $A$, $A'$, $A''$ conditions of SN-C04 are
unchanged.

In Grothendieck's later language, C44-F01 is a statement over
one scheme over $\mathbf Z$; (F03.2) is the comparison from
codimension-two cycles to \emph{\'etale motivic} cohomology;
(F02.2)--(F03.4) concern $\ell$-adic sheaves, weights, and their
arithmetic realization.  For the generic fibre, SGA~6,
Expos\'e~XIV, \S4.1--4.2 (especially (4.1) and (4.4)),
identifies the rational Chow group in codimension two with
degree two of the topological and $\lambda$ filtrations on
$K_0(J)$.  The operators in (F01.1) act there by pullback.
This $\lambda$-ring translation supplies no bound on rational
equivalences.  No new EGA/FGA theorem or standard-conjecture
implication is asserted: the modern comparison theorems are
cited under their actual sources above.

The open rows \texttt{QFermat-AJ} and \texttt{Q35-full} are
\emph{not closed}.  Their exact new test is (F03.4), or,
equivalently, bounded five-power denominators in $B$ on every
rational plane (F04.2).  The original analytic order-one result
AR-11, positive value AR-13, and Selmer-line result AR-23
remain intact.  None of the present findings proves
$\operatorname{rank}\operatorname{CH}^2(N)=1$.

% END work/arithmetic/contribution.tex

% BEGIN work/scan_w/contribution.tex
% C44_POS /root/scan_w.  New deductions, 2026-11-22 Europe/Berlin.
% The model is GPT-6 Astra Pro.  Runtime reasoning-effort setting is not exposed.
% This file is a paper fragment.  The root paper owns theorem environments.

\section{An algebraic replacement for Wirtinger in the Weil-plane detector argument}

The statements in this section distinguish an unconditional intersection
estimate from the positivity and algebraicity inputs.  Every invocation of
032, Corollaries 8.3 or 8.4, is conditional on the unrefereed preprints
specified by the user.  The characteristic-zero Theorem W in
\texttt{HODGE\_DIVISION\_ATTACK\_20261121.md}, R4, is not being reproved.
We give its missing positive-characteristic degree estimate and then record
the resulting, slightly differently normalized, detector theorem.

\subsection*{C44-W01. A characteristic-free intersection bound on an abelian variety}

\begin{proposition}[C44-W01]
Let $A$ be an abelian variety of dimension $2n$ over an algebraically
closed field, and let $h=c_1(L)$ be an integral ample class.  Put
\[
 D(Z)=\deg(h^n\cdot Z),\qquad
 V=\frac{\deg h^{2n}}{(2n)!}=\chi(L)\in\mathbf Z_{>0},\qquad
 b_h=\frac{2^n}{n!\sqrt V}.
\]
For effective integral cycles $Z,Z'$ of dimension $n$, including cycles
with multiplicities,
\begin{equation}\label{c44:algebraic-bezout}
 0\leq\deg(Z\cdot Z')\leq b_h^2D(Z)D(Z').
\end{equation}
No standard conjecture and no lifting hypothesis is used.
\end{proposition}

\begin{proof}
It suffices by bilinearity to treat integral subvarieties.  On
$Z\times Z'$ consider $f(z,z')=z-z'$ and $s(z,z')=z+z'$.  Translation
moving, or equivalently refined pullback of a general point of $A$ by
$f$, gives
\[
 f_*[Z\times Z']=e[A],\qquad e=\deg(Z\cdot Z')\geq0.
\]
If $f$ is not dominant, both sides have coefficient $e=0$.  If it is
dominant, it is generically finite and $e$ is its degree, including any
inseparable multiplicity.  To see explicitly that the intersection
coefficient agrees with this degree, a general translate of $Z'$ meets
$Z$ properly and avoids the singular loci of both subvarieties: the
bad incidence loci have dimension at most $2n-1$.  At the resulting
zero-dimensional intersections the two smooth subvarieties are local
complete intersections in the smooth variety $A$, and properness gives
the usual length intersection multiplicity.  The same description
applies to the generic fibre of $f$.  Translation preserves cycle
classes.  This also proves the assertion without separability or
transversality assumptions.

Let $h_1,h_2$ be the pullbacks of $h$ to $Z\times Z'$, and set
$\alpha=f^*h$, $\beta=s^*h$.  Both are nef.  The theorem of the cube,
in its parallelogram consequence, gives the numerical identity
\[
 \alpha+\beta=2(h_1+h_2).
\]
Here $[-1]^*h=h$ numerically; no symmetric representative of the
original line bundle is required.  A convenient exact source is
Milne, \emph{Abelian Varieties}, Corollary 6.5, applied to the maps
$z,z',-z'$; the numerical consequence is the displayed identity.
All mixed intersections of nef divisors with an effective cycle are
nonnegative, by approximation by ample rational classes.  Therefore
\begin{align*}
 e\deg h^{2n}
 &=\deg(\alpha^{2n}\cap[Z\times Z'])\\
 &\leq\deg((\alpha+\beta)^{2n}\cap[Z\times Z'])\\
 &=2^{2n}\binom{2n}{n}D(Z)D(Z').
\end{align*}
The last equality uses $\dim Z=\dim Z'=n$.  Division by
$\deg h^{2n}=(2n)!V$ proves the bound.  The asserted formula and
integrality for $V$ are the abelian Riemann--Roch theorem, Milne,
Theorem 13.3(b); ampleness gives $V>0$.
\end{proof}

\paragraph{Realizations and historical scope.}
Equation \eqref{c44:algebraic-bezout} is an inequality between rational
intersection numbers.  It is valid for classes represented in either
$\ell$-adic or crystalline cohomology because the cycle map preserves
these products and trace.  It does not put an order on $\mathbf Q_\ell$.
This is precisely the kind of degree bound needed below.  In
Grothendieck's framework its operations are proper pushforward, refined
intersection, and the biextension identity; SGA 7 I, Exp. VIII,
4.12.1 and 2.3.11 provide the biextension convention already used by
JAC-T4.3.  It supplies a new quantitative input to the positive
incidence calculation of cote 44, scans 008--009, and to Theorem W,
step 3.  It makes no assertion that a new axis class lifts.

\subsection*{C44-W02. The signed Weil-plane Gram matrix in any characteristic}

Fix an imaginary quadratic field $K$ and an order
$\mathcal O\subseteq K\cap\operatorname{End}(A)$ such that the Rosati
involution of $h$ induces conjugation on $K$.  Choose
$a\in\mathcal O$ with $a/\bar a$ not a root of unity.  Write
\[
 \nu=a\bar a\in\mathbf Z_{>0},\quad q=\nu^n,\quad
 s=a^{2n}+\bar a^{2n}\in\mathbf Z,\quad H=4q^2-s^2>0.
\]
For an invertible prime $\ell$, define
\[
 W_\ell=\bigwedge_{K\otimes\mathbf Q_\ell}^{2n}H^1(A,\mathbf Q_\ell)
 \subset H^{2n}(A,\mathbf Q_\ell).
\]
It has dimension two over $\mathbf Q_\ell$.  Let
\[
 f(t)=(t-a^{2n})(t-\bar a^{2n}),\quad
 g(t)=\prod_{j=1}^{2n-1}(t-a^j\bar a^{2n-j}).
\]
Both polynomials have integer coefficients.  Choose $R>0$ and
$u,v\in\mathbf Z[t]$ with $uf+vg=R$, for example the absolute
resultant with the sign of the Bezout relation adjusted.  Put
\[
 p(t)=v(t)g(t)=\sum_jc_jt^j,\qquad
 \kappa=\sum_j|c_j|\nu^{jn},\qquad e_W=R^{-1}p([a]^*).
\]
These are exactly the projector data of the supplied Theorem W.
We use them, with the same normalization, on the new realization.

\begin{proposition}[C44-W02; conditional on unrefereed 032, Corollary 8.4]
Let $Y\subset A$ be an integral $n$-dimensional subvariety with
$e_W[Y]\ne0$, and let $D=D(Y)$.  Then
\[
 y_1=p([a]^*)[Y],\qquad y_2=[a]^*y_1
\]
are independent integral algebraic cycle classes in $W_\ell$.
The form
\[
 Q(x,y)=(-1)^n\deg(xy)
\]
is positive definite on their rational span.  With $d=Q(y_1,y_1)$,
its Gram matrix is exactly
\begin{equation}\label{c44:weil-gram}
 G=d\begin{pmatrix}1&s/2\\s/2&q^2\end{pmatrix},\qquad
 \det G=d^2H/4\in\mathbf Z_{>0},\qquad
 0<d\leq b_h^2\kappa^2D^2.
\end{equation}
Moreover
\[
 W_{\rm alg}:=e_W\operatorname{Alg}^n(A)
    =\mathbf Qy_1\oplus\mathbf Qy_2,
 \qquad W_{\rm alg}\otimes\mathbf Q_\ell=W_\ell.
\]
The two classes have integral signed-cycle representatives whose
positive and negative parts have degrees at most $\kappa D$ and
$q\kappa D$, respectively.
\end{proposition}

\begin{proof}
After extending the coefficient field, write $H^1=V_a\oplus V_{\bar a}$.
Both eigenspaces have dimension $2n$, and the Rosati condition says
that $h$ lies in $V_a\otimes V_{\bar a}$.  Each of the two extreme
exterior powers is therefore annihilated by $h$.  Thus $W_\ell$ is
middle primitive.  The eigenspaces of $[a]^*$ in middle cohomology
have eigenvalues $a^j\bar a^{2n-j}$; all are distinct by the choice
of $a$.  Hence the integral polynomial $p$ is $R$ on the extreme
two lines and zero on the other weight spaces.  This verifies the
projector on this realization, rather than importing a Betti lattice.
The exterior-algebra description is valid integrally for etale
cohomology by Milne, Theorem 15.1.

Every $[a^j]$ is a finite flat isogeny, so the pullbacks in
$p([a]^*)[Y]$ are effective integral cycles with their flat
multiplicities.  Rosati gives $[a^j]^*h=\nu^jh$ numerically and
$\deg[a^j]=\nu^{2nj}$.  Projection formula consequently gives
\[
 D([a^j]^*Y)=\nu^{jn}D.
\]
This proves both asserted degree budgets.  The minimal polynomial
of $[a]^*$ on $W_\ell$ is the irreducible quadratic $f$: indeed
$a^{2n}$ is nonreal.  Hence $y_1,y_2$ are rationally independent.

The signed Hodge standard conjecture, used with exactly the sign
$(-1)^n$, is positive definite on every rational subspace of
primitive algebraic classes.  This is the sole input for positivity
here.  It is the even-degree version of the primitive sign in
cote 44, scan 008; JAC-T4.3 establishes the corresponding incidence
polarization by a curve test.  The unsigned form printed in
Milne's report, Section 3.6, is not used.

Let $T=[a]^*|W_\ell$.  Projection formula gives
$Q(Tx,Ty)=q^2Q(x,y)$.  Thus $T^\dagger=q^2T^{-1}$ and
$T+T^\dagger=s$.  It follows that
$Q(y_1,y_2)=sd/2$ and $Q(y_2,y_2)=q^2d$, proving the matrix
formula.  Its entries are integers, since the $y_i$ are integral
cycle classes, and its determinant is a positive integer.
Expanding $y_1^2$, applying C44-W01 to each pair of effective
pullbacks, and taking absolute values gives
\[
 d=|\deg(y_1^2)|
 \leq b_h^2D^2\Bigl(\sum_j|c_j|\nu^{jn}\Bigr)^2.
\]
This proves the quantitative replacement for Wirtinger.

Finally any finite rationally independent collection of classes in
$W_{\rm alg}$ has a positive definite rational Gram matrix by the
same Hodge-standard input.  Its determinant is nonzero, so those
classes are also $\mathbf Q_\ell$-independent.  There can be at most
two.  The two classes already constructed therefore span
$W_{\rm alg}$ and its realization.
\end{proof}

\paragraph{What ``integral'' means in positive characteristic.}
There is no canonical integral Betti lattice over $\overline{\mathbf F}_p$.
Let $\Lambda$ denote the group of integral algebraic middle classes
modulo numerical equivalence, viewed in cohomology using 032's
numerical--homological equality.  It is a lattice: choose finitely
many integral test cycles giving a nondegenerate rational pairing;
the intersection map embeds $\Lambda$ in a finite-rank free abelian
group.  Write
\[
 \Lambda^\sharp=\{x\in\operatorname{Alg}^n(A):\deg(x\lambda)\in\mathbf Z
                       \text{ for all }\lambda\in\Lambda\}.
\]
The denominator condition needed below is $\delta t\in\Lambda^\sharp$.
The stronger condition $\delta t\in\Lambda$ suffices.  If an integral
compatible realization is specified, its rational classes with integral
pairings also satisfy this condition.  Integrality only at one prime
$\ell$ is insufficient to make the Cramer numerators integers.

\subsection*{C44-W03. A quadratic detector bound in characteristic $p$}

\begin{theorem}[C44-W03; conditional on unrefereed 032, Corollary 8.4]
Retain the preceding notation and suppose $W_{\rm alg}\ne0$.
For $0\ne t\in W_{\rm alg}$ let $\delta\geq1$ be an integer with
$\delta t\in\Lambda^\sharp$, and put $T_0=Q(t,t)>0$.
Define $c_t$ by the same integral effective-cycle class equation as
in the supplied Theorem W, and let $\delta_W$ be its detector degree.
Set $B=b_h\kappa$.  Then
\begin{equation}\label{c44:weil-quadratic}
 \delta_W\leq c_t\leq
 \delta\max\{HB^2,\,4q\kappa B\sqrt{HT_0}\}\,\delta_W^2.
\end{equation}
In particular one can replace $B$ everywhere by the uniform constant
$B_0=2^n\kappa/n!$, independent of the abelian variety, the
characteristic, and the polarization type.  The quartic version
\begin{equation}\label{c44:weil-quartic}
 c_t\leq\delta q^2\max\{B^4,\,4\kappa B^3\sqrt{T_0}\}\,\delta_W^4
\end{equation}
also holds.  Thus the characteristic-$p$ division theorem has both
the requested detector formulation and a quadratic sharpening,
with fully algebraic, explicit constants.
\end{theorem}

\begin{proof}
The hypothesis $W_{\rm alg}\ne0$ supplies a detector: express a
nonzero class as a rational combination of integral effective
cycles and apply $e_W$.  Choose a detector $Y$ of minimal degree
$D=\delta_W$, and form $y_1,y_2$ as in C44-W02.  Put
$b_i=Q(\delta t,y_i)\in\mathbf Z$.  Inverting the exact Gram matrix
\eqref{c44:weil-gram} gives the identity
\begin{equation}\label{c44:exact-clearing}
 \delta Hd\,t=(4q^2b_1-2sb_2)y_1+(4b_2-2sb_1)y_2.
\end{equation}
All coefficients and the positive multiplier $N=\delta Hd$ are
integers.  If $t=t_1y_1+t_2y_2$, Cauchy--Schwarz in this
positive definite rational plane, or the diagonal entries of
$G^{-1}$, gives
\[
 |t_1|\leq\frac{2q\sqrt{T_0}}{\sqrt{dH}},\qquad
 |t_2|\leq\frac{2\sqrt{T_0}}{\sqrt{dH}}.
\]
Therefore the two integral coefficients on the right of
\eqref{c44:exact-clearing} have absolute values at most
$2\delta q\sqrt{HT_0d}$ and $2\delta\sqrt{HT_0d}$.
Using the degree budgets of C44-W02, the positive and negative
parts of the resulting integral cycle each have degree at most
\[
 \kappa D\bigl(2\delta q\sqrt{HT_0d}
                     +2\delta q\sqrt{HT_0d}\bigr)
 \leq4\delta q\kappa B\sqrt{HT_0}\,D^2.
\]
Also $N\leq\delta HB^2D^2$.  These two estimates prove the
upper bound in \eqref{c44:weil-quadratic}.

For completeness the ordinary two-by-two adjugate clearing with
multiplier $\delta\det G$ gives
$N\leq\delta q^2B^4D^4$.  Its two coefficients are bounded by
$2\delta\sqrt{T_0}\|y_1\|\|y_2\|^2$ and
$2\delta\sqrt{T_0}\|y_1\|^2\|y_2\|$, where
$\|y_1\|=\sqrt d\leq BD$ and $\|y_2\|=q\sqrt d$.
Their degree budget is at most
$4\delta q^2\kappa B^3\sqrt{T_0}D^4$, proving
\eqref{c44:weil-quartic}.

Conversely any class equation
$Nt=[Z^+]-[Z^-]$, with $N\geq1$ and effective integral cycles,
has a nonzero Weil projection on at least one integral component
of $Z^+$ or $Z^-$.  That component has degree at most the
corresponding cycle degree.  Hence $\delta_W\leq c_t$.
\end{proof}

\paragraph{Exact change to the supplied Theorem W.}
The original complex constant $|R|/n!$ in the norm estimate is
replaced by $2^n\kappa/(n!\sqrt V)$ for the integral numerator
$y_1$.  We do not claim that Wirtinger's sharper constant survives
unchanged.  All other projector, degree, independence, and integer
clearing steps are now valid without an archimedean realization.
The bound is on cohomological or numerical cycle classes, exactly
as in the supplied definition of $c_t$; it is not an equality in
the Chow group modulo rational equivalence.

\subsection*{C44-W04. The precise balanced-signature hypothesis over a finite field}

\begin{proposition}[C44-W04; Tate-algebraicity conclusion conditional on unrefereed 032]
Let $(A,h,K)$ descend to a finite field $\mathbf F_Q$, with the
$K$-action and polarization defined there, and $\dim A=2n$.
If $p$ splits in $K$, assume the two rational $K\otimes\mathbf Q_p$
components of $A[p^\infty]$ both have dimension $n$ (each has
height $2n$).  If $p$ is inert or ramified, no additional determinant
balance is necessary.  Then Frobenius acts with finite order on
$W_\ell(n)$.  Conditional on 032, Corollary 8.3, the whole
$W_\ell$ is algebraic, $W_{\rm alg}$ is its two-dimensional
rational cycle structure, $\delta_W<\infty$, and C44-W03 applies
to every nonzero class in that rational Weil plane.

When $p$ is split and $p$ does not divide the conductor of
$\mathcal O$, the stated $p$-divisible balance is exactly the usual
$(n,n)$ multiplicity condition on $\operatorname{Lie}A$.
It also holds for the special fibre of a balanced polarized
$\mathcal O_K\otimes\mathbf Z_p$-linear abelian scheme in mixed
characteristic.  These hypotheses give the usual meaningful
all-characteristic version of the signature condition.
\end{proposition}

\begin{proof}
Write $\pi$ for the $Q$-power Frobenius endomorphism.  Its
$K$-linear determinant on $H^1$ is an element $\alpha\in K$,
independent of the realization.  One can see its rationality
without choosing a complex lift: the traces of $b\pi^j$ for
$b\in K$ are rational, by the characteristic polynomial theorem
for endomorphisms, and the nondegenerate trace pairing of $K$
recovers the $K$-linear traces.  Newton identities recover the
$K$-linear characteristic polynomial and its determinant.
The same argument in crystalline cohomology gives the same
determinant: in any of these realizations, the exterior-algebra
description of abelian cohomology and the top-degree action give
$\det(m-[b])=\deg([m]-[b])$ for every integer $m$ for which the
map is an isogeny.  Polynomial identity then makes the
characteristic polynomial independent of the realization.
The trace/characteristic-polynomial input is
Milne, \emph{Abelian Varieties}, Proposition 12.9; it is the
degree-defined characteristic polynomial of an endomorphism.

The determinant $\alpha$ is an algebraic integer, since its
eigenvalues are products of Frobenius eigenvalues.  The
polarization pairing gives
\[
 \alpha\bar\alpha=Q^{2n},\qquad |\alpha|=Q^n
\]
at the two complex embeddings.  The absolute-value assertion is
Weil's theorem for abelian varieties, Milne, Theorem 19.1(c).
In particular $\alpha$ is a unit away from primes over $p$.
If $p$ is nonsplit, its unique prime is fixed by conjugation, so
the norm identity implies
$v(\alpha)=n\,v(Q)$ there.  If $p$ splits, the determinant of
Frobenius on each rational Dieudonne summand has normalized
determinant valuation equal to that summand's dimension:
the Newton polygon of a $p$-divisible group of height $h_0$
and dimension $d_0$ ends at $(h_0,d_0)$, so its slopes sum
to $d_0$.  For exact source statements see Oort,
\emph{Abelian Varieties over Finite Fields}, Sections
21.16--21.19 and Corollary 9.3.  Equivalently each
Dieudonne--Manin factor of dimension $d_0$ and height
$d_0+c_0$ has $d_0+c_0$ slopes $d_0/(d_0+c_0)$.
The balance hypothesis therefore gives
$v(\alpha)=n\,v(Q)$ at both primes.

Consequently $\alpha/Q^n$ is a unit in the imaginary quadratic
field $K$ of absolute value one, hence a root of unity.  On the
two Weil lines the Frobenius eigenvalues are $\alpha$ and
$\bar\alpha$.  A finite extension of the finite field kills
their normalized root-of-unity factors.  The Tate conjecture
in the stated 032 input then gives algebraic classes spanning
$W_\ell$: apply the algebraic projector $e_W$ to cycles spanning
the Tate classes.  C44-W02 identifies their rational cycle
span as a plane.  This proves the assertions.

For split $p$ and a $p$-maximal order, the two integral
idempotents act on the $p$-divisible group and on its Lie
algebra, so the dimensions agree with the two Lie
multiplicities.  In an integral balanced family these ranks
are constant.  Isogenies preserve the dimensions and heights
of the rational $p$-divisible summands.
\end{proof}

\begin{proposition}[C44-W05: a bad-order ambiguity in the word signature]
A determinant identity on $\operatorname{Lie}A$ for an order whose
conductor is divisible by a split prime $p$ does not by itself
imply the balance hypothesis of C44-W04.  In particular the
literal naive such identity cannot replace that hypothesis.
\end{proposition}

\begin{proof}
For a completely concrete example take
$E:y^2=x^3-x$ over $\mathbf F_5$, $K=\mathbf Q(i)$.
It has eight rational points, hence Frobenius trace $-2$;
in particular it is ordinary.  Its order-four automorphism
gives the $\mathbf Z[i]$-action, and its Frobenius is one
of $-1\pm2i$.  Thus $p=5$ splits in $K$, as required.
The following construction also works with any ordinary
CM elliptic curve at a split prime.  After a finite
extension all data are defined over $\mathbf F_Q$.
Let $A=E^{2n}$ with the scalar $K$-action and product polarization,
and restrict the action to $\mathcal O=\mathbf Z+p\mathcal O_K$
(or to $\mathbf Z+p\operatorname{End}(E)$ if necessary).
For $a=m+pb$ the action on the Lie algebra is multiplication by
$m\bmod p$, and the reductions of $a$ and $\bar a$ agree.
Thus
\[
 \det(T-a\mid\operatorname{Lie}A)
 =(T-a)^n(T-\bar a)^n
\]
holds after reduction for every $a\in\mathcal O$, although the
two rational $p$-divisible components have dimensions $2n$ and
$0$.  If $\pi$ is the ordinary Frobenius of $E$, the normalized
Weil determinant is $\pi^{2n}/Q^n$; its valuations at the two
primes over $p$ are nonzero and opposite.  It is not a root of
unity, and the two Weil lines are not Tate.  The example
therefore proves the claimed insufficiency.  It is not a
counterexample to C44-W03 or to the usual balanced PEL meaning
of Weil type.
\end{proof}

\section{Geometric positive transfers and special complete intersections}

The abstract abelian-type transfer remains distinct from a transfer that
preserves the actual incidence form. The following statements give
such geometric transfers and identify the extra special-fibre directions
that they do not reach. Their positivity conclusions are conditional
on the unrefereed 032 input whenever that qualification is displayed.

\subsection*{C44-W06. A finite abelian cover gives an explicit inherited polarization}

\begin{theorem}[C44-W06; conditional on unrefereed 032, Corollary 8.4]
Let $X$ be smooth projective of odd dimension $d=2m-1$ over an
algebraically closed field. Suppose that there is a finite surjective
morphism $f:A\to X$ from an abelian variety. For every rational ample
class $h$ on $X$, put $h_A=f^*h$ and $e=\deg f$. The primitive
incidence quotient $B_X=P^m(X)_{\rm prim}$ has the signed rational
polarization of cote 44, Th.~4. More precisely the induced map
\[
 i_f:B_X\longrightarrow B_A=P^m(A)_{\rm prim}
\]
has finite kernel, and its actual incidence classes satisfy
\begin{equation}\label{c44:finite-cover-polarization}
 i_f^*\ell_{B_A}=e\ell_{B_X},\qquad
 \ell_{B_X}=\frac{(-1)^m}{2}\Delta^*d_{B_X}.
\end{equation}
In particular this is a positive transfer on a concrete class of
nonabelian varieties of geometric abelian type. It does not use
Milne's additional transfer assertion from JAC-P20.
\end{theorem}

\begin{proof}
A finite morphism between smooth varieties of the same dimension is
flat; its degree includes inseparable multiplicity. Push--pull gives
$f_*f^*=e$, and $h_A$ is ample. Lefschetz inverses for $X$ are
algebraic: in degree $i\leq d$, an inverse to $L_h^{d-i}$ is
\[
 \frac1e f_* (L_{h_A}^{d-i})^{-1}f^*.
\]
Indeed $f^*$ commutes with the indicated Lefschetz operators and
the formula composes with $L_h^{d-i}$ to the identity. The inverse
on the abelian variety is algebraic. The same pull--push construction
with its Kunneth projectors supplies those of $X$. Thus the middle
primitive projector is algebraic and self-transpose. Saito's
incidence construction and transpose compatibility, used in
JAC-T4.1, give its nondegenerate factor $B_X$.

Pullback takes primitive cycle families on $X$ to primitive cycle
families on $A$, since $h_Af^*=f^*h$. It therefore induces $i_f$.
Pushforward induces a rational left inverse up to $e$, so $i_f$ has
finite kernel. For two normalized families $z,z'$ on parameter
varieties, projection formula gives the equality of divisorial
incidence correspondences
\[
 (p_{12})_*\bigl((1\times f)^*z\cdot(1\times f)^*z'\bigr)
   =e\,(p_{12})_*(z\cdot z').
\]
By the incidence functoriality of JAC-T4.1, this is the equality
$(i_f\times i_f)^*d_{B_A}=e d_{B_X}$. Diagonal pullback gives
\eqref{c44:finite-cover-polarization}, with the same sign and
normalization as JAC-T4.3. That theorem makes $\ell_{B_A}$ ample
conditional on 032. Pullback by a homomorphism with finite kernel
is ample on its source, and division by $e>0$ preserves rational
ampleness. This proves the conclusion. At the motive level
$e^{-1}f_*f^*=1$ exhibits $h(X)$ as an actual Chow-motive summand
of $h(A)$, so the stated abelian-type scope is also explicit.
\end{proof}

\paragraph{Concrete nonabelian examples.}
In characteristic different from two, take
$A=E_0\times E_1\times\cdots\times E_{d-1}$ with $d\geq3$ odd,
choose a nonzero two-torsion point $a\in E_0$, and quotient by
\[
 (x_0,x_1,x_2,\ldots)\longmapsto(x_0+a,-x_1,x_2,\ldots).
\]
The action is free because of its first coordinate. Its quotient
is smooth projective, the cover is finite etale, and the invariant
$H^1(\mathcal O)$ has dimension $d-1$, so the quotient is not an
abelian variety. In characteristic two use a nonzero three-torsion
translation on $E_0$ and an order-three automorphism of
$E_1:y^2+y=x^3$, acting by $x\mapsto\zeta_3x$. The same argument
works. These examples prove that C44-W06 has nonabelian content in
every characteristic.

\subsection*{C44-W07. The fixed part for every smooth complete intersection}

\begin{proposition}[C44-W07; fixed-part positivity conditional on unrefereed 032]
Let $A$ have dimension $N=2m+c-1$, let $h$ be ample, and let
$k:Y\hookrightarrow A$ be a smooth successive intersection of
$c$ divisors of class $h$, with $\dim Y=n=2m-1$. There is no
generic-last-section hypothesis in this statement. Define
\[
 r=(L_h^c)^{-1}k_*,\qquad
 F=k^*\ker L_h^{c+1},\qquad E=\ker k_*\subset H^n(Y),
\]
and let $e_A$ be the primitive projector on $H^n(A)$. Then
\[
 \Pi_F=k^*e_A r,\quad \Pi_E=1-k^*r,\qquad
 H^n(Y)_{\rm prim}=F\mathbin\perp E
\]
are algebraic self-adjoint projectors and an orthogonal decomposition.
They induce a rational isogeny decomposition of primitive incidence
quotients
\[
 P^m(Y)_{\rm prim}\sim B_F\times B_E.
\]
Its signed incidence class is nondegenerate on both factors and is
a rational polarization on $B_F$, conditional on 032. Consequently
the full Th.~4 assertion for this particular $Y$ is equivalent to
ampleness of its signed incidence class on $B_E$. In particular it
holds whenever $B_E=0$.
\end{proposition}

\begin{proof}
The fixed-and-variable projector proof of H9-T13 uses smooth weak
Lefschetz and the identity $[Y]=h^c$, before any pencil monodromy is
introduced. These hypotheses hold here. Explicitly,
$rk^*=1$, and $k^*r$ is the orthogonal projector onto the ambient
middle cohomology. The primitive condition for an ambient class
$k^*u$ is $L_h^{c+1}u=0$: the Gysin map on $H^{n+2}(Y)$ is
injective by dual weak Lefschetz. The same injectivity makes every
element of $\ker k_*$ primitive. The primitive projector of $A$
is algebraic; the projectors of $Y$ outside its middle degree come
from the weak-Lefschetz identifications with $A$ and their duals.
Thus the displayed middle projectors are actual algebraic
correspondences. Self-adjointness and nondegeneracy give the
incidence decomposition exactly as in JAC-T4.1.

For clarity, the positive scalar in the fixed-part test also remains
unchanged. If a normalized fixed family pulled to a pointed curve
$C$ of genus $b\geq1$ has class $\xi$, put $\eta=(1\otimes r)\xi$
and lift its parameter degree-one component to an algebraic class
$w$ on $J(C)\times A$. Then $h^{c+1}w=0$ and, for
$H=\theta+h$,
\[
 H^{b+c}w=0,\qquad
 \int_{J(C)\times A}w^2H^{b+c-1}
   =\frac{(b+c-1)!}{c!}\int_{C\times Y}\xi^2.
\]
This is H9-T13, equation (H9.47); its proof needs no genericity.
The coefficient is positive. Conditional on 032, the signed curve
degree is nonnegative, with equality only for a zero mixed class.
The nondegenerate signed divisor is therefore ample by the
nef-plus-nondegenerate criterion of JAC-T4.3. Orthogonality makes
ampleness of the full signed divisor equivalent to ampleness of
both factors, proving the final assertions.
\end{proof}

\subsection*{C44-W08. Specialization of the inherited incidence polarization}

\begin{theorem}[C44-W08; conditional on unrefereed 032]
Let $\mathscr Y\to U$ be the smooth locus of a genuine Lefschetz
pencil in the setting of H9-T13, and let $s\in U$ be any geometric
point. Write $B_{\bar\eta}=P^m(Y_{\bar\eta})_{\rm prim}$ for the
geometric generic primitive incidence quotient. After a finite
extension of the local generic field, $B_{\bar\eta}$ has good
reduction at $s$, and its special fibre maps with finite kernel to
$P^m(Y_s)_{\rm prim}$. The actual signed incidence form is positive
on this inherited image. Hence Th.~4 holds for $Y_s$ whenever
\begin{equation}\label{c44:no-incidence-jump}
 \dim P^m(Y_s)_{\rm prim}=\dim P^m(Y_{\bar\eta})_{\rm prim}.
\end{equation}
More particularly, if the generic variable incidence factor of the
pencil is nonzero, then Th.~4 holds for every smooth geometric fibre
of the pencil. If that generic factor is zero, only the new
variable incidence directions that occur at a special fibre remain
outside this theorem.
\end{theorem}

\begin{proof}
Work over the strict henselian local trait at $s$, or its completion,
and pass to a finite extension of its fraction field over which an
abelian representative $B$ and a normalized rational Poincare cycle
$z$ are defined. These are finite data; rational descent is the one
fixed in JAC-P00. As in H9-T5, equations (H9.10)--(H9.11), the
mixed class of $z$ injects the Tate module of $B$ into the appropriate
Tate twist of $H^n(Y_{\bar\eta})$. Indeed its self-incidence is
nondegenerate, so a vector killed by the family map would be in
the radical of that incidence pairing. Smooth proper base change
makes the target unramified at this smooth fibre. Thus $B$ has
good reduction by the Neron--Ogg--Shafarevich criterion, just as in
H9-T5. One exact primary author source for the criterion is Milne,
\emph{Complex Multiplication}, Theorem 6.12, p.~51, which cites
Serre--Tate (1968), Theorem 1.

Let $\mathscr B$ denote the resulting abelian scheme over the trait.
Take the closure of a rational-cycle representative of $z$ in
$\mathscr B\times\mathscr Y$. These schemes are smooth over a
regular trait, so the refined fibre restrictions and intersection
products are defined. Its special-fibre mixed class is the
specialization of the generic mixed class. In particular it is
still in parameter degree one and is primitive on $Y_s$; these
conditions are kernels of maps of the smooth-proper cohomology
systems. Rigidification or degree-one projection, if desired,
alters no incidence class.

Form the actual relative self-incidence divisor by the same proper
pushforward of the product of the two cycle families. Clear its
rational denominator. Its signed half-diagonal pullback is a line
bundle $\mathscr L$ on $\mathscr B$, and refined base change says
that its two fibre classes are exactly the corresponding signed
incidence classes, up to the one positive clearing factor. The
generic class is ample by H9-T13.

Here is the specialization step without an assertion that arbitrary
special cycles lift. Choose a relatively ample line bundle
$\mathscr H$ on $\mathscr B$. Flatness and the Hilbert polynomial
make
\[
 \chi(L_t+uH_t)=\frac{(c_1(L_t)+u c_1(H_t))^{\dim B}}{(\dim B)!}
\]
the same real polynomial in $u$ for the generic and special fibres
$t$. On the generic fibre this polynomial is positive for every
$u\geq0$, since $L_{\eta}+uH_{\eta}$ is ample. Thus on the
special fibre $\chi(L_s)\ne0$ and the polynomial has no positive
real root. Mumford's index theorem gives index zero for $L_s$;
the nondegenerate effective criterion makes it ample. These are
exactly the two abelian-variety criteria used in the checked proof
of JAC-T4.3 (Mumford, \emph{Abelian Varieties}, Chapter III,
Section 16, and Chapter II, Section 6, Application 1).

The specialized family induces a homomorphism
$u:B_s\to P^m(Y_s)_{\rm prim}$. Saito's transpose and incidence
compatibility identifies $u^*\ell_s$ with the ample class just
constructed. Therefore $u$ has finite kernel: restriction to any
positive-dimensional connected kernel would make an ample class
zero. This proves positivity on the inherited image. Under
\eqref{c44:no-incidence-jump} the map is an isogeny, and ampleness
descends through that isogeny.

For the final assertion apply the same construction to the generic
variable factor alone. If that factor is nonzero, the irreducibility
after finite-index restriction in H9-T13, or H9-T5 (H9.11), forces
its family map to fill the entire variable local system $E$.
Specialization preserves this rank. The special variable incidence
factor injects into $E_s$ by its own nondegenerate self-incidence,
so its dimension is at most $\operatorname{rank}(E)/2$. The
inherited image already has that dimension. It therefore maps
isogenously onto the special variable incidence factor and gives
its positive signed form. C44-W07 supplies the fixed factor for
every smooth fibre. This proves the stated all-smooth-fibre scope.
\end{proof}

\paragraph{The remaining special-fibre gate.}
When the generic variable incidence quotient vanishes, the
construction specializes only the fixed incidence quotient. A
special fibre may have a new variable incidence factor. Neither
smooth proper base change nor constancy of the rank of $E$ makes
that factor extend as algebraic cycle families through the pencil.
Its signed positivity is the smaller unresolved statement. Thus
C44-W08 does not turn H9-T13 into an unconditional assertion about
every smooth complete intersection; it proves precisely the
rank-preserving specializations and the nonzero-generic-variable
case, with 032 retained as the positivity input.

\subsection*{C44-W09. The full geometric-abelian-type extension with its extra input exposed}

\begin{proposition}[C44-W09; conditional on unrefereed 032 and the separate JAC-P20 transfer]
Assume, in addition to 032, Milne's Proposition 4.2 in
\emph{Grothendieck's standard conjecture of Lefschetz type over
finite fields}, version 14 July 2022, in the total-cohomology form
specified and audited in JAC-P20. Then JAC-T4.3 holds for every
smooth projective odd-dimensional variety of geometric abelian
type, with its actual signed primitive incidence divisor.
\end{proposition}

\begin{proof}
JAC-P16 gives algebraic Lefschetz and primitive projectors. For
every pointed test curve $C$, the product $J(C)\times X$ is still
of geometric abelian type. JAC-P20's actual abelian cohomology
injection and the additional transfer assertion give the Hodge
standard conjecture for this product. The construction of
JAC-T4.2, using the algebraic Abel inverse on the first factor,
now gives an algebraic primitive class $w$ on $J(C)\times X$ and
the exact identity
\[
 \int_{J(C)\times X}w^2(\theta+h)^{b-1}
       =(b-1)!\int_{C\times X}\xi^2.
\]
The Hodge-standard input on this particular product makes its
signed value nonnegative. Saito's nondegenerate incidence factor
and JAC-T4.3's curve criterion then make the signed divisor ample.
This argument proves the stated conditional corollary; it does
not establish the additional transfer assertion itself.
\end{proof}

\paragraph{Framework and precise dependency.}
C44-W06 proves a positive scalar comparison of the actual forms,
not merely an embedding of motives. C44-W07--W08 use smooth
proper base change, SGA 7's fixed and variable projectors, and
the extension of actual cycle correspondences across a regular
trait. Their target is cote 44, scans 008--009, Th.~4 and its
primitive incidence higher Picard variety. The new residual gate
is the positive sign on a special new variable incidence factor,
not the already checked generic pencil sign. For all varieties
of geometric abelian type, JAC-P20 still leaves the comparison
between an abstract transported positive form and the particular
geometric form as an additional dependency. C44-W09 records
the exact implication if that dependency is supplied.


\section{The supersingular Fermat test for cote 16}

\subsection*{C44-S01. The proposed Fermat candidate already has the required sign}

\begin{proposition}[C44-S01; unconditional]
Let $X=X_N^{2r}\subset\mathbf P^{2r+1}$ be the smooth Fermat
hypersurface of degree $N\geq3$ in characteristic $p\nmid N$,
with $r\geq2$.  Suppose $p^\nu\equiv-1\pmod N$ for some
$\nu\geq1$.  Then all even cohomology is algebraic and all odd
cohomology vanishes.  Numerical and homological equivalence
coincide.  The signed form
\[
 Q_V(x,y)=(-1)^r\deg(xy),\qquad
 V=\ker(h:H^{2r}(X)\longrightarrow H^{2r+2}(X)),
\]
is positive definite on its rational algebraic cycle structure.
In particular the Hodge standard conjecture for $X$ is known
unconditionally; it is not a consequence that needs 032.
\end{proposition}

\begin{proof}
Put $q=p^\nu$ and $b=(q+1)/N$.  The Hermitian Fermat
hypersurface $H=X_{q+1}^{2r}$ maps finitely and surjectively to
$X$ by
\[
 [x_0:\cdots:x_{2r+1}]
 \longmapsto[x_0^b:\cdots:x_{2r+1}^b].
\]
Its degree is $b^{2r+1}$: $f^*h_X=bh_H$ and
$\int_Hh_H^{2r}=q+1=bN$ determine the degree by projection
formula.  In particular $f_*f^*=b^{2r+1}$ and $f^*$ is injective.
Shimada, \emph{Lattices of Algebraic Cycles on Fermat Varieties
in Positive Characteristics}, Theorem 1.1(1)--(2) and
Corollary 1.2, proves that linear cycles span the middle
cohomology of $H$ and that its signed primitive intersection
form is positive definite.  Weak Lefschetz gives the other
cohomology groups.  Pushing forward these algebraic classes
shows algebraicity of all cohomology of $X$.

If $x\in V$, then
$h_Hf^*x=b^{-1}f^*(h_Xx)=0$, and
\[
 (-1)^r\int_H(f^*x)^2
  =b^{2r+1}(-1)^r\int_Xx^2.
\]
The injectivity of $f^*$ and Shimada's positive form prove the
claim.  In all other degrees $X$ has the ambient projective
space Lefschetz strings, whose sign is immediate.  The
intersection pairing is nondegenerate on the rational
algebraic classes because these span cohomology; thus
numerical and homological equivalence agree.  Since
$\operatorname{NS}(X)_\mathbf Q=\mathbf Qh$, this also covers
every ample polarization.
\end{proof}

\paragraph{Exact status of the proposed test.}
Shioda--Katsura, \emph{On Fermat Varieties}, Theorem II,
printed p.~97, says for even dimension and $N\geq4$ that
algebraic generation of cohomology is equivalent to the
displayed congruence.  Thus, in that range, the entire stated
supersingular Fermat candidate class is covered by C44-S01.
The new argument above is only the finite-cover descent
of Shimada's already known sign.  It neither attributes a
new theorem to Grothendieck nor treats algebraicity alone as
a positivity theorem.  The square needed in scan 049 is
handled explicitly below.

\subsection*{C44-S02. An explicit smooth Segre hyperplane}

\begin{proposition}[C44-S02; unconditional]
Let $X=X_N^d$ be a smooth Fermat hypersurface with $N\geq3$.
Choose a finite field $k_0$ containing every $(N-2)$-nd root
of unity in $\overline{\mathbf F}_p$.  Let $M>N(d+1)$ and
choose $\beta$ of degree $M$ over $k_0$ in its extension of
degree $M$.  Define $c_i=\beta^{(N-2)i}$ for
$0\leq i\leq d+1$.  Then
\[
 W=\left\{\sum_{i=0}^{d+1}c_ix_iy_i=0\right\}
       \subset X\times X
\]
is a smooth hyperplane section for the Segre polarization
$L=h_1+h_2$.  Also $Y=X\cap\{x_{d+1}=0\}$ is a smooth
hyperplane section of $X$.
\end{proposition}

\begin{proof}
The assertion for $Y$ follows from $p\nmid N$.  If $W$ had
a singular point, the differential of its equation would
annihilate both projective tangent spaces.  Since the
Fermat gradients are $N(x_i^{N-1})$ and $N(y_i^{N-1})$,
there would be scalars $\lambda,\mu$ with
\[
 c_iy_i=\lambda x_i^{N-1},\qquad
 c_ix_i=\mu y_i^{N-1}\quad\text{for every }i.
\]
Both scalars are nonzero, since all $c_i$ are nonzero and
neither projective coordinate vector is zero.  The nonzero
coordinates of $x$ and $y$ have the same nonempty support
$S$.  Substitution gives, for $i\in S$,
\[
 (x_i^N)^{N-2}=(\mu\lambda^{N-1})^{-1}c_i^N.
\]
Choose $\rho$ with $\rho^{N-2}=(\mu\lambda^{N-1})^{-1}$.
It follows that
$x_i^N=\rho\zeta_i\beta^{Ni}$ for roots of unity
$\zeta_i\in k_0$ with $\zeta_i^{N-2}=1$.
This remains valid if $p\mid(N-2)$, since only the distinct
roots are used.  The Fermat equation for $x$ now gives
\[
 \sum_{i\in S}\zeta_i\beta^{Ni}=0.
\]
This is a nonzero polynomial over $k_0$ of degree at most
$N(d+1)<M$ vanishing at $\beta$, a contradiction.  The
Jacobian criterion therefore proves smoothness.
\end{proof}

\subsection*{C44-S03. The exact square sign and all three corrected scan-050 conditions}

\begin{theorem}[C44-S03; unconditional]
For $X=X_N^{2r}$ of C44-S01, put $d=2r$, $T=X\times X$,
$L=h_1+h_2$, and take $Y,W$ as in C44-S02.  Then
\[
 C_{\rm G}(X),\quad C_{\rm G}(Y),\quad R_0(T,W),\quad
 G_{L,0}(T,W),\quad A_0(T),\quad G_0(T,W),\quad A_0(W)
\]
all hold with the meanings fixed by SIGN\_NEXT SN-C03--SN-C04.
The Hodge standard conjecture $I(T)$ holds.  In its middle
degree the primitive space and form are explicitly
\begin{gather}
 P^{2d}(T)=\mathbf Qw\ \perp\ (V\otimes V),\qquad
 w=\sum_{i=0}^{d}(-1)^ih_1^ih_2^{d-i},\label{c44:fermat-square-primitive}\\
 \int_T(cw+z)^2=(d+1)N^2c^2+(Q_V\otimes Q_V)(z,z)>0
 \quad\text{if }(c,z)\ne(0,0).\label{c44:fermat-square-sign}
\end{gather}
The ring of degree-zero rational correspondence classes
modulo homological equivalence is
\[
 \prod_{i=0}^{d}\operatorname{End}_{\mathbf Q}
       \bigl(\operatorname{Alg}^i(X)\bigr)
 \simeq \mathbf Q^d\times M_{b_d(X)}(\mathbf Q),
\]
and is therefore semisimple.  These are the actual sign and
semisimplicity conclusions of cote 16, scan 049, for this
nonabelian family.
\end{theorem}

\begin{proof}
Weak Lefschetz gives the orthogonal decomposition
\[
 H^{\rm even}(X)=U\perp V,\qquad
 U=\mathbf Q\langle1,h,\ldots,h^d\rangle,
\]
where $V$ is in cohomological degree $d$ and $hV=0$.
The rational structures here are given by cycle classes;
they become the full cohomology after scalar extension.
The middle cohomology of $T$ consequently consists of the
ambient terms, the two terms $V\otimes h^r$ and
$h^r\otimes V$, and $V\otimes V$.  On the ambient
terms the kernel of multiplication by $L$ is the single
line displayed in \eqref{c44:fermat-square-primitive}:
successive coefficients satisfy $c_i+c_{i-1}=0$.
Multiplication by $L$ is injective on each of the two
cross terms and their images have different Kunneth
degrees.  It is zero on $V\otimes V$.  Orthogonality
follows from $\int_Xvh^r=0$ for $v\in V$.

In $w^2$ the $d+1$ contributing pairs have indices
$i,d-i$ and sign $(-1)^d=1$, giving
$\int_Tw^2=(d+1)N^2$.  The pairing on $V\otimes V$
is the tensor product of the two raw pairings; their
two factors $(-1)^r$ cancel.  It is therefore
$Q_V\otimes Q_V$, a positive definite form.  This proves
the exact critical sign.

For the full Hodge standard statement, decompose
$H^*(T)$ into the four $L$-stable blocks
$U\otimes U$, $V\otimes U$, $U\otimes V$, and
$V\otimes V$.  They are orthogonal.  The first has
the Lefschetz module and intersection forms of
$\mathbf P^d\times\mathbf P^d$, multiplied by $N^2$;
its signed forms are positive by characteristic-zero
Hodge--Riemann, with identical rational matrices in
every characteristic.  In $V\otimes U$ the primitive
classes are exactly the base $V\otimes1$, in
codimension $r$.  Their signed form is $NQ_V$;
the same is true for $U\otimes V$.  The last block
is primitive in middle degree and was just treated.
This proves $I(T)$ in all degrees.  Every ample
class on $T$ is a positive linear combination of
$h_1,h_2$; the same argument, with positive rescalings,
also handles these polarizations if required.

For a smooth hypersurface $Z$ of degree $N$ and
dimension $e$, the Lefschetz inverses below its middle
degree are already explicit.  For even $i<e$ the
correspondence
\[
 \frac1N h_Z^{i/2}\times h_Z^{i/2}
\]
induces the inverse of
$L_Z^{e-i}:H^i(Z)\longrightarrow H^{2e-i}(Z)$;
for odd $i<e$ both spaces vanish.  The middle inverse
is the identity.  Thus $C_{\rm G}(X)=B(X)$ and
$C_{\rm G}(Y)=B(Y)$.  The product stability of $B$
(Kleiman 1968, Proposition 2.12), or the cycle-basis
matrix units below, gives $B(T)$ and in particular
the algebraicity of $(L^2)^{-1}$ in the critical
degrees of SN-C03.

Write $j:W\hookrightarrow T$.  If $j^*x$ is algebraic
for $x\in H^{2d-2}(T)$, then
$L j_*j^*x=L^2x$ is algebraic; applying the algebraic
inverse proves $R_0(T,W)$.  For
$\gamma\in\operatorname{Alg}^{d+1}(T)$, set
$x=(L^2)^{-1}\gamma$ and $y=j^*x$.  Both are
algebraic and $Lj_*y=\gamma$, proving $G_{L,0}$.
The critical $A_0(T)$ was used here; the equivalence
in SN-C03 supplies $G_0$ and $A_0(W)$ as well.
This verifies the maps with the required extra $L$;
no unweighted lower-Gysin replacement is made.

Finally choose rational cycle bases for each
$\operatorname{Alg}^i(X)$ and the intersection-dual
bases in degree $d-i$.  Exterior products of these
classes give all rational matrix units on
$H^{2i}(X)$.  Every degree-zero correspondence
acts degreewise; Kunneth and Poincare duality make
this action faithful modulo homological equivalence.
It is therefore exactly the product of the stated
full matrix algebras.  Since all nonmiddle Betti
numbers in even degree are one, the displayed
description follows.
\end{proof}

\paragraph{What task 4 does and does not close.}
C44-S01--C44-S03 rule out the suggested supersingular Fermat
family as an example whose square-sign status is unknown.
They verify the corrected scan-050 proposition and the
scan-049 sign and correspondence ring explicitly and
unconditionally on that family.  A different smooth
projective $X$ with a genuinely unestablished $I(X\times X)$
is still required for the adversarial part of task 4; absence
of a source for one named equation would not certify that
status.  None of these findings closes the special-axis
sign for a complete intersection in an abelian variety.

\paragraph{Framework and programme crosswalk.}
For C44-S01--S03 the relevant archive locators are
\texttt{ARCHIVE\_16\_scan049.tex}, the semisimple
correspondence paragraph, and
\texttt{ARCHIVE\_16\_scan050.tex}, the proposition and
final weighted diagram.  They instantiate SIGN\_NEXT
SN-C03--SN-C04 and its successor question 4.  Their
Lefschetz inverses also give the $C_{\rm G}$ input to
cote 44, Th.~6, scan 004, through JAC-P02.  Weak
Lefschetz and Gysin are the Weil-cohomology operations
in SGA 5, Exp. VII, Section 7 (complete intersections).
The finite power maps and explicit hyperplanes are
scheme morphisms over finite extensions of $\mathbf F_p$;
the identities survive every Weil cohomology realization.
At the level of homological motives the Fermat object
is a direct sum of Tate objects with the positive
Lefschetz forms just computed.  No equality with its
Chow motive modulo rational equivalence is silently
deduced.  C44-W01--W05 use the exterior-power, hence
$\lambda$-operation, of $h^1(A)$, the algebraic
endomorphism projector, and rational cycle intersections.
No additional $K$-theory assertion, Selmer injection,
or motivic $L$-function nonvanishing is inferred.

% END work/scan_w/contribution.tex

\part{Unchecked candidates and unresolved work}
The following two cover candidates are awaiting an independent whole-proof reading. They are not included among checked findings at this checkpoint.
% C44_POS /root; candidates undergoing independent whole-proof review.
\subsection{C44-C03: characteristic-two conductor strata with a simple pole}
\begin{theorem}[C44-C03]\label{thm:AS-strata}
In characteristic two, fix an Artin--Schreier double-cover pole partition
\[
 d_i=2a_i+1\quad(1\le i\le s),\qquad
 s\ge2,\quad a_i\ge0,\quad a_1=0,
 \qquad g=\sum_i(a_i+1)-1\ge2.
\]
Fix an abelian-fourfold Lefschetz pencil as in H9-T22, and require the
branch locus of the covers to avoid its critical values. At the geometric
generic point of the indicated irreducible conductor stratum, there is
no algebraic class in the new middle trace kernel of the pulled-back
pencil. Consequently
\[
 h_L(z,w)=\tfrac12h_K(\operatorname{cor}z,\operatorname{cor}w),
 \qquad \operatorname{rad}h_L=
 \operatorname{cor}^{-1}(\operatorname{rad}h_K).
 \tag{C44.C03.1}
\]
The identity and vanishing are unconditional. Semipositivity of the
resulting height uses H9-T8 and is conditional on the unrefereed
032/Milne inputs. The assertion includes the geometric generic point
of the hyperelliptic $2$-rank $g-1$ divisor for every $g\ge2$.
It is a statement at geometric generic points, not at every closed
point of the stratum.
\end{theorem}
\begin{proof}
Order the poles and use the dense coordinate chart in which all poles
are finite, with the distinguished simple pole at $t=0$. Put the other
poles at distinct nonzero $b_2,\ldots,b_s$, away from the critical set
$\Sigma$ of the pencil. Choose a rational function $f_0$, regular at zero
and infinity, with exactly these poles and orders $d_2,\ldots,d_s$.
The family
\[
                  w^2-w=f_0(t)+\lambda/t,
                  \qquad \lambda\in\mathbf G_m
 \tag{C44.C03.2}
\]
lies in the specified conductor stratum. Its normalization is smooth
and proper over the parameter. Here is the required local verification
also for a higher pole. If $z$ is a local parameter at a pole of order
$2a+1$, set $v=z^{a+1}w$. The integral local equation becomes
\[
 v^2-z^{a+1}v=z^{2a+2}f(z).
\]
The right side has first term $c z$, with $c$ a unit on the parameter
space. At the unique point $z=v=0$ its derivative in $z$ is $c\ne0$.
The equation is regular there, finite flat of degree two, and agrees
with the extension off the pole, so it is the normalization. Away from
the poles the derivative in $w$ is a unit. This proves the relative
smoothness and justifies smooth proper base change.

Let $V$ be the variable odd local system of the original pencil,
of rank $r>0$, with $q=\#\Sigma$. Its irreducibility and rank-one tame
transvections are the inputs used in H9-T21--T22. Change coordinates
to $x=1/t$, and put
\[
 G=\mathcal L_\psi(f_0(1/x)),\qquad
 F=(\operatorname{inv}^*V)\otimes G.
\]
Their middle extensions to $\mathbf A^1_x$ are irreducible Fourier
sheaves: $G$ has a finite wild pole, and $F$ is irreducible of rank
$r\ge2$. Both are unramified at infinity. The exact character-part
Leray calculation of H9-T22 identifies the new middle cohomology
with
\[
 H^3(A)\otimes\operatorname{NFT}_\psi(G)
       \ \oplus\ \operatorname{NFT}_\psi(F)
 \tag{C44.C03.3}
\]
over the parameter open. This is a calculation of the actual
smooth-proper cohomology, not a substitution of extension by zero at
the old transvections. The outer Leray terms are deck invariant; the
two displayed terms are its nontrivial deck character.

Set $R=\sum_{i=2}^s(d_i+1)$. The Fourier ranks are $R$ and $q+rR$,
by the finite conductor formula, Katz, \emph{Exponential Sums and
Differential Equations}, Lemma 7.3.9(2). Theorem 7.3.8(3) gives
irreducibility of the transforms. At parameter zero, Corollary
7.4.3.1 and Theorem 7.4.4(2)--(3), equivalently Theorem 7.5.4(5)--(7),
give their respective local representations
\[
 J(1,2)\oplus J(1,1)^{R-2},\qquad
 J(1,2)^{r}\oplus J(1,1)^{q+rR-2r}.
 \tag{C44.C03.4}
\]
Each is infinite because it contains a nontrivial unipotent block.
An irreducible representation with infinite image has no nonzero
finite-orbit vector: the span of the orbit would be a nonzero
subrepresentation, and the normal core of its finite-index stabilizer
would then act trivially on the whole representation. This would make
its image finite. Direct sums and tensoring with a constant vector
space preserve the absence of finite-orbit vectors. Thus (C44.C03.3)
has none.

For completeness, the space with ordered pole positions and all
reduced odd principal-part coefficients is irreducible: it is an open
in a vector bundle over the ordered configuration space, with each
leading coefficient nonzero. The above local normalization works
throughout it. It maps dominantly to the indicated conductor stratum
of covers. Its universal trace-kernel local system restricts to
(C44.C03.3). A finite-orbit vector for the universal monodromy would
remain one on this parameter curve. There are therefore no such
vectors at the geometric generic point of the stratum. Algebraic
classes, defined over finite extensions of its function field, have
finite monodromy orbit. The actual trace projector and admissibility
argument of H9-T22 and H9-T24 now give (C44.C03.1).

The genus and $2$-rank are $g$ and $s-1$. Pries--Zhu,
\emph{The $p$-rank stratification of Artin--Schreier curves},
Theorem 1.1 and Corollary 1.3, printed p.~708, identify the components
in characteristic two with partitions of $g+1$ into $s$ positive
parts $a_i+1$, all of dimension $g-1+(s-1)$. For rank $g-1$ the
partition is $(2,1,\ldots,1)$, with $g-1$ parts equal to one.
It has a simple pole and is a divisor in the hyperelliptic moduli
space of dimension $2g-1$. Fixing the target adds three dimensions
to both spaces and leaves codimension unchanged. This proves the
last assertion.
\end{proof}

\subsection{C44-C04: extremal triple-cover splitting strata}
\begin{theorem}[C44-C04]\label{thm:extremal-maroni}
Let $e\ge2$ and $g=3e-2$. In every characteristic, let $U_{e,2e}$
denote the framed space of smooth connected triple covers of the
fixed target $\mathbf P^1$ with Tschirnhausen bundle
$\mathcal O(e)\oplus\mathcal O(2e)$, unramified over the old
critical set. At its geometric generic point the new middle
trace kernel of an abelian-fourfold pencil contains no algebraic
classes. The norm-height identity and radical in (C44.C02.1) hold.
These assertions are unconditional; height semipositivity is
conditional on the unrefereed 032/Milne inputs through H9-T8.

Together with C44-C02 this settles the geometric generic point
of every admissible triple-cover splitting stratum of genus at
least two. It does not assert the same vanishing at every closed
point over $\overline{\mathbf F}_p$.
\end{theorem}
\begin{proof}
The cubic coefficient space for $E=\mathcal O(e)\oplus\mathcal O(2e)$
is $H^0(\mathcal O\oplus\mathcal O(e)\oplus\mathcal O(2e)
\oplus\mathcal O(3e))$. Its smooth connected unramified locus is
irreducible if nonempty. We construct a family in it for each
characteristic and prove that its new trace-kernel monodromy has
no finite-orbit vectors. The restriction argument of C44-C02
will then prove the assertion for this whole stratum.

\paragraph{Characteristic different from three.}
Choose a coordinate $t$ with zero and infinity outside the old critical
set $\Sigma$. Choose $3e-3$ nonzero points $b_i$ so that their
$\mu_3$-orbits are pairwise disjoint and avoid the $\mu_3$-orbits of
$\Sigma$. Put $f(t)=\prod_{i=1}^{3e-3}(t-b_i)$ and consider
\[
                   u^3=(t^3-s)f(t).
 \tag{C44.C04.1}
\]
Remove zero and the finitely many parameters for which a moving root
hits a fixed branch point or a critical value. The normalized covers
are smooth and have $3e$ simple branch points, all with exponent one;
infinity is unramified. Their genus is $3e-2$, and the relative
eigensheaf decomposition is
\[
        \pi_*\mathcal O_C=
        \mathcal O\oplus\mathcal O(-e)\oplus\mathcal O(-2e).
 \tag{C44.C04.2}
\]
The normalized local equations are the tame equations of H9-T23.
Because the degree of the branch polynomial is $3e$, this is a
relative decomposition with fixed line bundles, not only a
fibrewise numerical calculation. It supplies a framing into $U_{e,2e}$.

Fix a nontrivial cubic character $\chi$. On the $t$ line let
$G_\chi=\mathcal L_\chi(f(t))$ and
$F_\chi=V\otimes G_\chi$, always taking middle extensions.
For $\rho:\mathbf A^1_t\to\mathbf A^1_v$, $v=t^3$, set
\[
                  M_\chi=\rho_*G_\chi,
                  N_\chi=\rho_*F_\chi.
\]
Both are irreducible middle extensions. Indeed, restrict to the
$\mu_3$-stable open on which $\rho$ is a Galois cover. Each input
is irreducible. Its three conjugates are pairwise nonisomorphic:
at a fixed branch point one has nontrivial scalar inertia $\chi$,
whereas either conjugate is lisse, by the choice of the $b_i$.
The elementary induction criterion for a normal subgroup with
pairwise inequivalent conjugates therefore makes its induced
representation irreducible. Finite pushforward is its middle
extension, as is checked by inertia invariants at the omitted points.

Finite pushforward and projection formula identify the two new
character terms with
\[
          H^3(A)\otimes\operatorname{MC}_\chi(M_\chi)
                    \ \oplus\ 
                    \operatorname{MC}_\chi(N_\chi).
 \tag{C44.C04.3}
\]
This follows from the actual H9-T23 character cohomology by replacing
$\mathcal L_\chi(t-s)$ there with
$\mathcal L_\chi(t^3-s)$ and then applying $\rho_*$. The fibre
cohomology is the middle convolution, including its local invariant
conditions. Both input sheaves are tame. At infinity their
representations are respectively
\[
                 1\oplus\chi\oplus\chi^2,
                (1\oplus\chi\oplus\chi^2)^{r}.
 \tag{C44.C04.4}
\]
Here the original $V$ and $G_\chi$ are unramified at infinity,
and the ramification of $t\mapsto t^3$ supplies the regular
representation of $\mu_3$.

Katz, \emph{Rigid Local Systems}, Theorem 3.3.3 and Corollary 3.3.6,
give irreducibility of the convolutions and their local invariant
quotient formula. In that formula write
$M_\infty/M_\infty^{I_\infty}$ for the input representation.
Each trivial one-dimensional summand in (C44.C04.4) forces a
$J(1,2)$ block in $M_\infty$. Tensoring by the nontrivial $\chi$
turns it into $J(\chi,2)$, and quotienting by inertia invariants
does not remove it. Thus both outputs have infinite local image.
For a rank check, the first input has finite drops $2$ at zero
and $1$ at each of the $3e-3$ fixed branch images; Corollary
3.3.7 gives output rank $3e-2$. The second has finite drops
$q+(3e-1)r$ in total and output rank $q+(3e-2)r$.
Collisions among old critical images do not change their total
drop; the chosen branch images are disjoint from them.
Irreducibility and infinite image exclude finite-orbit vectors,
as in C44-C03.

\paragraph{Characteristic three.}
Choose $e-1$ distinct nonzero points $b_j$, away from $\Sigma$, so
that the reciprocal points $b_j^{-1}$ and their negatives are
distinct and avoid the reciprocal critical set and its negatives.
For nonzero constants $c_j$, put
\[
 f_0(t)=\sum_{j=1}^{e-1}\frac{c_j}{(t-b_j)^2},
 \qquad w^3-w=f_0(t)+s/t^2,\qquad s\in\mathbf G_m.
 \tag{C44.C04.5}
\]
There are $e$ poles, all of order two, hence genus $3e-2$.
The local order-two normalization and the global basis calculation
of H9-T23--T24, (H9.131) and (H9.144), give a smooth family and
the relative splitting (C44.C04.2). Its pole divisor is fixed,
so once again it maps to the framed stratum.

Put $x=1/t$. For either nonzero additive character index
$\alpha\in\mathbf F_3^\times$, set
\[
 G_\alpha=\mathcal L_\psi(\alpha f_0(1/x)),\qquad
 F_\alpha=(\operatorname{inv}^*V)\otimes G_\alpha,
 \qquad \rho:x\longmapsto v=x^2.
\]
The two finite pushforwards $\rho_*G_\alpha$ and
$\rho_*F_\alpha$ are irreducible by the same induction argument,
now for the tame quadratic cover and its involution $x\mapsto-x$.
A finite wild pole at $b_j^{-1}$ has no pole in its negative orbit,
so the two conjugates cannot be isomorphic. Their ranks are $2$
and $2r$, so neither is an exceptional rank-one additive character.
They are Fourier sheaves. At infinity their representations are
$1\oplus\epsilon$ and $(1\oplus\epsilon)^r$, respectively,
where $\epsilon$ is the nontrivial quadratic tame character.

Projection formula converts the variable term
$\mathcal L_\psi(\alpha s x^2)$ into
$\mathcal L_\psi(\alpha s v)$. The actual new character
cohomology is consequently the sum of $H^3(A)$ tensored with
$\operatorname{NFT}_\psi(\rho_*G_\alpha)$ and
$\operatorname{NFT}_\psi(\rho_*F_\alpha)$, pulled back by
$s\mapsto\alpha s$. Katz, \emph{Exponential Sums and Differential
Equations}, Theorem 7.3.8(3), gives irreducibility. Corollary
7.4.3.1 and Theorem 7.4.4(2)--(3) show that the trivial infinity
summands produce $J(1,2)$ blocks at Fourier parameter zero.
The nontrivial quadratic summands remain nontrivial rank-one
blocks and do not cancel these unipotent blocks. Hence these
transforms also have infinite image and no finite-orbit vectors.

In both characteristics every new character has now been handled.
The direct sum has no finite-orbit vectors. The universal lisse
trace-kernel projector on $U_{e,2e}$ restricts to this sum, so
the restriction argument proves generic algebraic vanishing on
the stratum. The norm-height identity and radical follow by
the same actual-cycle push-pull as in C44-C02. All smooth connected
triple-cover splittings have $e\le f\le2e$; the inequality follows
also directly from the cubic equation, since if $f>2e$ its first
coefficient is zero and the cubic has a linear factor, contradicting
smooth connectedness. C44-C02 handles the strict inequality, while
the present result handles equality for all genus at least two.
\end{proof}

\paragraph{Programme location and residual cover question.}
C44-C03--C04 extend the monodromy method of JAC1964 H9-T21--T24 to
specified special strata. Their height meaning is cote 44, Th.~5,
scan 004, with the primitive qualification of scan 008; the
idempotent $1-b^{-1}g^*g_*$ is the new Chow-motive projector.
The argument concerns lisse sheaves and finite monodromy orbits,
using SGA~7~II, XVIII~6.7.2 for the original pencil variation as
in the predecessor proofs. It gives no positivity theorem for an
extra Tate class at a closed point. Such a class belongs to the
remaining H9-T18 primitive incidence gate, after C44-C01 removes
the positive old surface-axis summand.

\section{Consolidation under final review}
\section{The corrected fifth theorem, with all domains visible}

\begin{theorem}[C44-T05: synthesis of H9-T1, H9-T2, H9-T7,
H9-T17 and H9-T18]\label{thm:final-th5}
Use rational Chow groups and the signed geometric intersection height
of JAC1964 H9-T1. The following is the consolidated form of cote 44,
scan 004, Th.~5. The formulas below are unconditional under their
stated admissibility and geometric hypotheses. Each assertion that
uses the abelian Hodge standard conjecture is explicitly conditional
on the unrefereed 032/Milne inputs.

\begin{enumerate}
\item Let $f:\mathcal X\to C$ be projective with $\mathcal X$ smooth
of dimension $2m$, $C$ a smooth projective curve, and smooth generic
fibre $Y$. Let $x_0,y_0$ extend numerically trivial cycles on $Y$.
For each relevant fibre let $W_s$ be its vertical numerical cycle
space, $R_s$ its radical, and $M_s$ the nonsingular matrix on a chosen
basis of $W_s/R_s$. Suppose the extensions annihilate $R_s$, and
that these spaces account for all vertical pairings. Write $a_s,b_s$
for their pairing vectors. Then admissible corrections exist and
\[
 h_m(x_0|_Y,y_0|_Y)=(-1)^m
 \left(\deg(x_0y_0)-\sum_s a_s^{t}M_s^{-1}b_s\right).
 \tag{C44.T05.1}
\]
If only a smaller vertical space has been used, this expression is
the corrected intersection, and is a height only after the remaining
vertical orthogonality is verified. Failure to annihilate a vertical
radical is an obstruction to correction, not a term one may omit.

\item For a Lefschetz pencil on a smooth $2m$-fold $X$, if
$H^{2m}(Y_s)\to H^{2m}(Y_{\bar\eta_s})$ is injective at every
singular fibre, all corrections vanish for ambient cycles whose
generic restrictions are homologically trivial:
\[
                  h_m(j^*z,j^*z')=(-1)^m\int_Xzz'.
 \tag{C44.T05.2}
\]
The same injectivity makes any global extension of a homologically
trivial generic cycle admissible. In the ordinary Picard--Lefschetz
setting the condition is equivalent to nonzero vanishing cycles;
SGA 7 condition $(A)$ is the characteristic-independent formulation
used in the predecessor. For genuine abelian pencils it holds by
JAC-P07. Conditional on 032/Milne, (C44.T05.2) is positive definite
on primitive ambient algebraic cohomology.

\item If $X$ is abelian, let $B$ be its smooth axis and use
$\mathcal H,\mathcal H_0,A_X,A_B$ of H9-T7, (H9.17). Then
\[
 \mathcal H/\mathcal H_0\simeq A_X\mathbin\perp A_B,
 \qquad h_m\simeq
 \left((-1)^m\int_Xzz'\right)\mathbin\perp
 \left((-1)^{m-1}\int_Bww'\right).
 \tag{C44.T05.3}
\]
In particular $h_Bw=0$ for $w\in A_B$. Conditional on 032/Milne,
semipositivity on the whole $\mathcal H$ is equivalent to the weak
axis sign on $A_B$. If the axis form is positive definite, the
radical is exactly $\mathcal H_0$. If it is only semidefinite,
the quotient of the radical by $\mathcal H_0$ is the radical of
that axis form; equality with $\mathcal H_0$ is not inferred.
The axis sign is known for $m=2$, for characteristic zero, and in
the generic-axis vanishing cases H9-T8, H9-T8a and JAC-P09.
C44-A04 and C44-A06 give additional explicit positive subspaces,
not a sign theorem for every special axis.

\item For a smooth threefold generic fibre, on the admissible
homologically trivial domain let
\[
 \mathcal H_h^{\rm Alb}=\{z:\operatorname{alb}(hz)=0\},\qquad
 \psi_h(D)=\operatorname{alb}(h^2D),\qquad
 D_z=\psi_h^{-1}\operatorname{alb}(hz).
\]
The rational isogeny $\psi_h:\Pic^0(Y)\to\Alb(Y)$ gives the
orthogonal decomposition
\[
                    z=(z-hD_z)+hD_z.
 \tag{C44.T05.4}
\]
The first term lies in $\mathcal H_h^{\rm Alb}$; the second has
negative definite height modulo the constant trace. In
characteristic zero the first height is semidefinite by Zhang,
Theorem 1, with the precise admissible-model equality criterion
of H9-T17, (H9.77)--(H9.80). This sign statement does not extend
to the whole homologically trivial group after an arbitrary
finite cover: H9-T17 gives the exact negative value $-12d^3$ on
the displayed $E^4$ example, in every characteristic.

\item For an abelian-fourfold pencil and a finite separable cover
$C\to\mathbf P^1$ unramified over the critical values, let
$k:\mathcal X\hookrightarrow A\times C$ be the smooth incidence
fourfold. All homologically trivial generic cycles are admissible.
On $\mathcal H_h^{\rm Alb}$ the primitive correction of H9-T18
gives
\[
              h_2(z,z)=\int_{A\times C}a^2\lambda
                              +\int_{\mathcal X}v^2,
 \quad \lambda^2a=0,\quad k_*v=0.
 \tag{C44.T05.5}
\]
Conditional on 032/Milne the first term is positive definite on
its algebraic domain. The remaining assertion is precisely the
nonnegativity of the algebraic primitive form on $\ker k_*$.
C44-C01 splits off its positive old surface-axis part; the new
trace-kernel part is the remaining special-cover gate. The
generic vanishing results H9-T19--T26 and C44-C02--C04 give
the norm identities in their explicitly stated parameter scopes.
They do not establish this sign for every closed special cover.
\end{enumerate}
\end{theorem}

\begin{proof}
This is an assembly of the specified predecessor results, not a new
proof of them. Part (1) is H9-T1, (H9.3)--(H9.5), equal to
SIGN\_NEXT SN-44-01 after its normalization, with vertical blocks
separated by disjoint fibres. Part (2) is H9-T2, (H9.6)--(H9.7),
and its full-extension argument at the beginning of H9-T7;
the Picard--Lefschetz equivalence is SN-44-03. Part (3) is the
isometry H9-T7, (H9.18), with the elementary radical statement
of an orthogonal sum and the conditional ambient sign. Part (4)
is H9-T17, (H9.79)--(H9.84), including both of Zhang's hypotheses
and Kahn's trace quotient. Part (5) is H9-T12 admissibility and
H9-T18, (H9.86)--(H9.91), followed by the independently checked
old/new decomposition C44-C01. No step asserts algebraic
equivalence from homological triviality.
\end{proof}

\paragraph{The incidence higher Picard variety and the arithmetic reading.}
Grothendieck's $P^m$ in cote 44, scans 002 and 004, is attached
to algebraically trivial cycles via incidence. A homologically
trivial class in the displayed larger Chow groups is not thereby
a point of $P^m$. The Neron--Tate interpretation on a primitive
incidence quotient requires the actual regular homomorphism and
the height comparison; these are the separate inputs of
JAC-T4.1--T4.5 and SN-44-01. The full-middle-factor polarization
is refuted by JAC-P03, and the unrestricted covered height is
refuted by H9-T17. Neither counterexample refutes the primitive
axis sign asked here. The normalization and the Albanese kernel
implement the primitive-decomposition margin of scan 008.

In cote 16, scans 049--050, the algebraicity of Lefschetz operators
is condition $C=C_G=B$, while the height signs use $I$ on the
specified primitive algebraic spaces. Restoring the Lefschetz
factor in $A''$ (SN-C03--C04) is a separate issue from these
signs. Theorem C44-T05 does not deduce $I$ from $C$ or identify
cohomological semisimplicity with injectivity on rational Chow
groups. In the SGA framework, the local input is SGA 7 II,
XVIII 5.3.2, 5.3.4--5.3.5 and 6.2; the variable irreducibility
is XVIII 6.7.2. In the Chow-motive framework the global terms
are the blowup summand $h(B)(1)$ and the old/new finite-cover
projectors. These locators and their roles are retained from
the checked predecessor arguments.

\section{Unresolved goals}
The fixed-axis sign is not proved for every special axis, and no
wrong-sign algebraic class is constructed. The universal lifting
question for classes remains open. The special-cover sign remains
open at arbitrary closed points with new algebraic trace-kernel
classes. Fermat Abel--Jacobi injectivity remains open, with the
exact gates specified above. Full Th.~4 for arbitrary new special
incidence directions is not proved. The proposed supersingular
Fermat family does not meet the requested initially-unknown-$I$
status. The ledger and question register give the current precise
state of every treated row.
\end{document}
